% !TeX program = pdflatex
%==============================================================================
% A Chronology of the HELM Reproducibility Internship (Kitware, Summer 2026)
%
% Author:  Edward Wang <edward.wang@kitware.com>
% Purpose: An authoritative, detailed technical chronology of the internship
%          work, written as a series of technical reports / blog posts, to
%          serve as source material for an academic paper on benchmark
%          reproducibility.
%
% This document is reconstructed from: the weekly intern status decks
% (docs/internship-powerpoints.zip), the append-only developer journals
% (dev/journals/{claude,codex}.md and their 2026-H1 archives), the git history
% (452 authored commits, 2026-06-03 .. 2026-07-14), the standing research
% journal (docs/helm-reproduction-research-journal.md), and the technical
% design docs under docs/ (helm-gotchas.md, vllm-vs-huggingface-deployment-
% match.md, helm-unrecorded-deployment-params.md, eee-vs-helm-metadata.md,
% pipeline.md, architecture.md, container-execution.md).
%==============================================================================
\documentclass[11pt]{report}

\usepackage[margin=1in]{geometry}
\usepackage[T1]{fontenc}
\usepackage[utf8]{inputenc}
\usepackage{microtype}
\usepackage{xspace}
\usepackage{amsmath}
\usepackage{amssymb}
\usepackage{amsfonts}
\usepackage{booktabs}
\usepackage{tabularx}
\usepackage{longtable}
\usepackage{array}
\usepackage{float}
\usepackage{xcolor}
\usepackage{listings}
\usepackage{enumitem}
\usepackage{graphicx}
\usepackage[
  colorlinks=true,
  linkcolor=black,
  citecolor=black,
  urlcolor=blue
]{hyperref}
\usepackage{bookmark}
\usepackage{url}
\usepackage{xurl}

\newcolumntype{Y}{>{\raggedright\arraybackslash}X}
\newcolumntype{Z}{>{\raggedright\arraybackslash}p{0.30\textwidth}}

\lstdefinestyle{shell}{
  basicstyle=\ttfamily\scriptsize,
  breaklines=true,
  breakatwhitespace=false,
  columns=fullflexible,
  keepspaces=true,
  showstringspaces=false,
  frame=single,
}
\lstset{style=shell}

% Shorthands
\newcommand{\eee}{\textit{Every Eval Ever}\xspace}
\newcommand{\eeeshort}{\textsc{eee}\xspace}
\newcommand{\helm}{\textsc{helm}\xspace}
\newcommand{\code}[1]{\texttt{#1}}
\newcommand{\file}[1]{\texttt{#1}}
\newcommand{\fromspec}{\emph{from-spec}\xspace}

% A consistent "finding" callout
\newcommand{\finding}[1]{\par\noindent\textbf{Finding.}~#1\par}

\title{\bfseries Reproducing \textsc{helm}: A Technical Chronology\\[4pt]
  \large What Portions of a Public LLM Benchmark Are Independently\\
  Reproducible Under a Corrected Local Open-Weight Recipe,\\
  and Where Unavoidable Differences Appear\\[10pt]
  \normalsize Kitware Research Internship --- Summer 2026}
\author{Edward Wang\\ \texttt{edward.wang@kitware.com}\\[4pt]
  \small Kitware, Inc.}
\date{Chronology covering 2026-05-15 through 2026-07-14}

\begin{document}
\maketitle

\begin{abstract}
\noindent
This document is a detailed, chronological technical report of a summer research
internship whose goal evolved from ``perform large-scale reproductions of the
\textsc{helm} benchmark'' into a sharper scientific question: \emph{which portions
of the public \textsc{helm} corpus are independently reproducible under a corrected
local open-weight execution recipe, and where do unavoidable differences appear?}
The central methodological contribution is a disciplined separation between
\emph{recipe/environment failures} (a model is gated, a dataset is no longer
fetchable, an engine cannot serve a required precision) and \emph{true
reproducibility failures} (same recipe, same data, same model, different metric).
The central engineering contribution is a pipeline that replays a public run's
\file{run\_spec.json} \emph{verbatim} inside a version-pinned container, serves the
model on independently controlled local hardware, and quantifies the residual
disagreement with per-metric agreement curves, tolerance sweeps, and
aggregate-score-drift heatmaps. Along the way we catalogue a family of concrete,
publishable ``gotchas'' --- silent \file{run\_spec.json} schema evolution, an
unrecorded execution substrate (dtype, model revision, attention kernel,
chat-template version), a tokenizer that appends an end-of-text token and produces
prompt-independent gibberish, an official \textsc{olmo} run that was secretly
\code{float32}, a chat-template version drift keyed on \code{add\_generation\_prompt},
and a classic-corpus class path that resolves in no released \textsc{helm} version ---
each of which is exactly the kind of thing a reviewer of a reproducibility paper will
ask about. We also summarize a parallel, exploratory research thread on
\emph{efficient benchmarking} (feature selection plus regression, and a Data-Kernel
Perspective-Space baseline). The document is written as a series of self-contained
technical reports so that individual findings can be lifted directly into an
academic manuscript.
\end{abstract}

\tableofcontents

%==============================================================================
\chapter{How to Read This Document}
%==============================================================================

\section{Scope and provenance}

This is an \emph{authoritative reference} assembled after the fact from primary
sources so that a future paper can cite it with confidence. The reconstruction
draws on:

\begin{itemize}[nosep]
  \item the seven weekly intern status decks (2026-06-02 through 2026-07-14),
        which record the work in the author's own words and contain the headline
        result figures;
  \item the append-only developer journals (\file{dev/journals/claude.md} and its
        \file{2026-H1} archive), which capture, per work session, the intent, the
        design alternatives considered, the chosen approach, and the reusable
        lesson --- roughly one hundred design narratives;
  \item the git history: 452 commits authored between 2026-06-03 and 2026-07-14,
        whose messages give a fine-grained, dependency-ordered record of what
        landed and when;
  \item the standing research journal
        (\file{docs/helm-reproduction-research-journal.md}), which frames the
        scientific question and records the early foundational findings; and
  \item the technical design documents accumulated under \file{docs/}, chiefly
        \file{helm-gotchas.md}, \file{vllm-vs-huggingface-deployment-match.md},
        \file{helm-unrecorded-deployment-params.md}, \file{eee-vs-helm-metadata.md},
        \file{pipeline.md}, \file{architecture.md}, and
        \file{container-execution.md}.
\end{itemize}

The internship ran from 2026-05-15. The first $\sim$2.5 weeks (through the 06-02
status update) were spent on onboarding and a literature review; the first
authored commits appear on 2026-06-03. The work then proceeds essentially without
interruption through 2026-07-14, the last date this chronology covers.

\section{A note on inherited foundations versus internship contributions}

The internship did not start from an empty repository. A substantial foundation
was already in place: the \file{eval\_audit} package itself; the \emph{Every Eval
Ever} (\eeeshort) normalized artifact format used as the common comparison
substrate; the \code{kwdagger} scheduler integration; and an early set of
reproducibility findings on \textsc{boolq}/Pythia and Vicuna/\textsc{NarrativeQA}
(described in \S\ref{sec:inherited}). Where this document describes that
foundation, it does so to make the internship's own contributions legible; the
chronology proper (Chapters~\ref{chap:e2e} onward) is the author's work. The
developer journals are written in the first person by the AI pair-programming
harness (Claude Code) that the author drove; throughout this report, ``we'' denotes
the author directing that toolchain, making the design decisions, and executing and
interpreting the runs on the GPU hosts.

\section{The two threads}

Two research threads ran in parallel for the first month, after which the second
was paused in favor of the first:

\begin{description}[style=nextline,leftmargin=1.5em]
  \item[Thread A --- \textsc{helm} reproduction (Chapters \ref{chap:framing}--\ref{chap:reporting}).]
    The main line of work: building the machinery to faithfully replay public
    \textsc{helm} runs on local open-weight hardware, and using it to determine
    what is and is not reproducible, and why. This is the bulk of the internship
    and of this document.
  \item[Thread B --- efficient benchmarking (Chapter \ref{chap:effbench}).]
    An exploratory statistics thread: can a small, well-chosen \emph{coreset} of
    benchmark questions plus a regression predict a model's full-benchmark score,
    and does a Data-Kernel Perspective-Space (\textsc{dkps}) representation help?
    This lived in two sibling repositories (\file{mrmr\_eval}, \file{dkps}) and
    reached a working prototype with preliminary results before being paused.
\end{description}

%==============================================================================
\chapter{The Research Question and Its Framing}
\label{chap:framing}
%==============================================================================

\section{From ``reproduce \textsc{helm}'' to a falsifiable claim}

The stated initial goal was to perform large-scale reproductions of the public
\textsc{helm} (Holistic Evaluation of Language Models) results. The public
corpus publishes, for many (scenario, model) pairs, the prompts sent, the
completions returned, and the computed metrics. Reproducing a result means
re-running the same evaluation locally and checking that the numbers match.

That goal did not survive contact with the data, and the way it failed is itself
the scientific content. Two categories of models dominate the public corpus:
open-weight models that were originally served through a hosted API (frequently
\emph{Together AI}), and closed-weight models that cannot be run locally at all.
The recipe that produced each public number is only partially recorded: the
\file{run\_spec.json} captures \emph{what} to generate (prompts, decoding
parameters, scenario, metrics) and the \emph{name} of the model and its
deployment, but almost nothing about the \emph{execution substrate} --- the
precision the weights were loaded at, the exact checkpoint revision, the attention
kernel, the tokenizer and chat-template versions, the software-stack versions. As
a result, three distinct difficulties compound:

\begin{enumerate}[nosep]
  \item \textbf{Running the benchmarks at all} --- gated models and datasets,
        packaging incompatibilities, download endpoints that have since revoked
        access, and scenario code that requires optional dependencies or
        credentials.
  \item \textbf{Setting up a \emph{faithful} local reproduction} --- matching the
        original recipe when key execution parameters were never written down, and
        when the original was served by an external provider whose stack we cannot
        inspect.
  \item \textbf{Attributing disagreement} --- when local and public numbers
        differ, deciding whether the benchmark is legitimately non-reproducible,
        whether the difference is an artifact of an uncontrollable external
        deployment, or whether our local reproduction is simply missing a detail.
\end{enumerate}

The goal was therefore reframed. The paper-worthy claim is \emph{not} ``all of
public \textsc{helm} is reproducible.'' It is:

\begin{quote}\itshape
A carefully defined, runnable subset of \textsc{helm} appears reproducible under a
corrected local open-weight recipe, with clearly documented failure modes.
\end{quote}

\section{The load-bearing distinction: two kinds of failure}
\label{sec:distinction}

Everything downstream --- the Stage-1 filters, the Sankey diagrams, the
per-packet diagnoses, the aggregate heatmaps --- is organized around one
distinction that must be held rigorously.

\begin{description}[style=nextline,leftmargin=1.5em]
  \item[Recipe / environment failures (these are \emph{filtering reasons}, not
        reproducibility results).] The model is gated behind a license; the dataset
        is gated or its download endpoint denies access; the model is packaged in a
        way local deployment cannot consume (e.g.\ a vision--language model); the
        required GPU architecture or precision is unavailable; a scenario needs a
        proprietary annotator. A run excluded because of
        \code{no-local-helm-deployment} is a \emph{design constraint}, not evidence
        that \textsc{helm} is irreproducible.
  \item[True reproducibility failures.] Same recipe, same data, same model, but
        different metrics. These are measured with a tolerance sweep from
        \code{abs\_tol}$=0$ (exact match) out to \code{abs\_tol}$=0.1$ (a tenth of
        the metric's range), analyzed per metric to see which metrics are fragile,
        and validated across machines to separate platform-specific from
        recipe-specific drift. An \code{agreement\_ratio} that dips to $0.7$ at
        \code{abs\_tol}$=0$ is a genuine reproducibility problem.
\end{description}

Much of the intellectual work of the internship was correctly \emph{classifying}
each observed disagreement into one of these buckets --- and resisting the
temptation to fold every failed large-grid job into a blanket negative result.
Chapters~\ref{chap:deploymatch} and~\ref{chap:eras} are, in large part, the story
of moving specific disagreements out of the ``irreproducible'' bucket and into the
``environment/recipe'' bucket once their true cause was found.

\emph{Refinement (interpretive revision, Appendix~\ref{app:revision}).} This binary
is the right first cut, but it is too coarse for the paper, and the phrase ``true
reproducibility failure'' should be avoided wherever execution provenance is
incomplete. Keep the upstream run-vs-not-run gate (non-runnable is an eligibility
reason, not a disagreement), but resolve observed disagreements into a finer
six-category taxonomy --- recipe / deployment / execution-instrument /
artifact-migration mismatch, residual disagreement under controlled equivalence, and
\emph{historically non-identifiable reproduction}. Only the fifth resembles
conventional repeatability failure; the sixth is qualitatively different (the target
of ``faithful reproduction'' is underdetermined). See Appendix~\ref{app:revision}.

\section{The inherited baseline: what was already known}
\label{sec:inherited}

Before the internship, the project had established a small but load-bearing set of
results, recorded in \file{docs/helm-reproduction-research-journal.md}. They set
the priors that the internship's work refined.

\begin{itemize}
  \item \textbf{Local reruns are stable; the official--local gap is not noise.} For
        \code{boolq:model=eleutherai/pythia-6.9b}, two local repeats agreed at a
        strict run-level ratio of $0.955$ with a maximum run-level delta of
        $0.0015$, whereas official-vs-local agreement was $0.463$ with run-level
        deltas up to $\sim\!12$. The official--local gap is far larger than
        rerun noise, so ordinary nondeterminism cannot be the main explanation ---
        the residual is structural (deployment and execution-spec drift).
  \item \textbf{A dramatic ``irreproducibility'' turned out to be a local
        misconfiguration.} Local Vicuna/\textsc{NarrativeQA} produced 99/100 empty
        completions; the cause was \textsc{helm} silently auto-applying a
        \emph{chat template} to a non-chat prompt. Rerunning with
        \code{apply\_chat\_template: false} recovered near-official scores
        (\code{exact\_match} $0.27$, \code{f1} $0.64$, \code{rouge\_l} $0.64$).
        The lesson --- that \code{apply\_chat\_template} must be an explicit,
        controlled setting, and that high empty-completion rate / near-zero output
        tokens / pervasive \code{finish\_reason\_unknown} are red flags --- recurs
        throughout the internship.
  \item \textbf{The metric-scale caveat.} A single global tolerance is hard to
        interpret because \emph{core} metrics live in $[0,1]$ while
        \emph{bookkeeping} metrics (byte counts, token counts, log-probs) do not.
        The reporting must separate metric classes; this became the
        \code{metrics\_taxonomy} module (\S\ref{sec:taxonomy}).
\end{itemize}

These priors motivated the internship's dominant methodological choice --- prefer
open-weight models runnable on realistic hardware, control the deployment
explicitly rather than trusting \textsc{helm}'s automatic inference, and treat the
reproducibility question as a \emph{distribution} (same-machine, cross-machine,
official-vs-local) rather than a single point estimate.

%==============================================================================
\chapter{Onboarding and Literature Review (mid-May -- June 2)}
\label{chap:litreview}
%==============================================================================

The first status deck (2026-06-02) reports two workstreams: end-to-end tests for
the audit pipeline (Chapter~\ref{chap:e2e}) and a literature review that seeded
both the reproducibility thread and the efficient-benchmarking thread
(Chapter~\ref{chap:effbench}). Three papers were reviewed in depth.

\begin{description}[style=nextline,leftmargin=1.25em]
  \item[LoRA weight-space learning --- \emph{W2T: LoRA Weights Already Know What
        They Can Do} (Han et al., 2026), building on Chen, Helm, Park, Priebe \&
        Villar, \emph{Extracting information from fine-tuned weights}.]
    A LoRA update is $AB^\top$; the factorization $(A,B)$ is non-identifiable
    because $(AM, BM^{-\top})$ yields the same update for any invertible $M$ (a
    $\mathrm{GL}$-group symmetry). W2T canonicalizes via the SVD and embeds the
    canonicalized weights with a transformer into a ``weight-space,'' whose
    embeddings drive downstream tasks: attribute classification (properties of the
    fine-tuning data), performance prediction, and adapter retrieval. The review
    noted that the ablations attribute most of the method's power to the
    \emph{canonicalization} step, and floated directions: generative modelling over
    LoRA weights, simpler statistical embeddings (e.g.\ MDS) in place of the
    transformer, and multi-/merged-adapter settings.
  \item[Efficient benchmarking --- \emph{Efficient Benchmarking Is Just Feature
        Selection and Multiple Regression} (Bowyer, Locatelli \& Cao, 2026).]
    Given $M$ source models, represent each question $X_i$ as a vector in
    $\mathbb{R}^M$ (each coordinate is a model's score on that question) and the
    target $Y$ as the vector of overall scores; \emph{select} an informative
    subset of questions with mRMR (Minimum-Redundancy-Maximum-Relevance) and
    \emph{predict} full-benchmark scores with ridge / kernel-ridge regression.
    Baselines: Lasso, Anchor Points, IRT. This framing became Thread~B.
  \item[Connecting the two.] The review proposed comparing LoRA and \textsc{dkps}
        (Data-Kernel Perspective-Space) representations, using Bayesian / Gaussian-
        process regression to unify selection and regression, and enriching the
        per-question feature vector with question-text embeddings --- directions
        pursued in Chapter~\ref{chap:effbench}.
\end{description}

%==============================================================================
\chapter{End-to-End Tests: Validating the Audit Instrument}
\label{chap:e2e}
%==============================================================================

\emph{Commits: 2026-06-03 through 2026-06-04, and the E2E deck of 2026-06-02.}

\section{Motivation}

Before trusting the audit pipeline's verdicts, the pipeline itself had to be
validated end to end on a case where the ground truth was controllable. The
audit workflow has six stages: (1) stand up a vLLM server for the target model
(via \code{infer-stack}, over LiteLLM / docker-compose); (2) write a benchmark
bundle; (3) execute runs on the servers; (4) aggregate run artifacts; (5) index
them; (6) build the analysis. The tools involved --- \code{infer-stack},
\code{eval-audit}, \code{magnet}, \code{kwdagger}, \code{cmd\_queue} --- each have
seams where a subtle mismatch produces a plausible-but-wrong comparison.

\section{The controlled experiment}

The subject was \code{microsoft/phi-2} on \code{mmlu:philosophy}, run under three
deliberately chosen deployments:

\begin{itemize}[nosep]
  \item \textbf{HF}: \textsc{helm}'s in-process \code{HuggingFaceClient},
        \code{temperature=0};
  \item \textbf{vLLM}: served through vLLM, \code{temperature=0} --- should match
        HF up to engine numerics;
  \item \textbf{vLLM ``incomparable''}: served through vLLM at \code{temperature=1}
        --- a deliberate negative control that \emph{should} disagree, proving the
        instrument can detect a real difference.
\end{itemize}

\section{Three instrument bugs found and fixed}

Validating the instrument surfaced three ways the comparison machinery could
silently mispair runs --- each fixed so that equivalent runs are actually
recognized as equivalent:

\begin{enumerate}
  \item \textbf{Hard-coded serving ports.} \code{infer-stack} used hard-coded
        ports without exposing them in a form \textsc{helm} could consume. Fix:
        read ports from the configuration source of truth and export them into a
        \code{.env} file queryable by the rest of the pipeline.
  \item \textbf{The \code{--groups} flag.} Public benchmark runs carry a
        \code{--groups} flag that has no effect on execution but prevented
        equivalent local runs from being matched to their public counterparts. Fix:
        ignore \code{groups=} tokens when locating comparable public runs, and emit
        a warning that this normalization occurred.
  \item \textbf{The \code{--temperature} run-name elision.} \textsc{helm} strips
        the temperature from the run \emph{name}, creating a divergence between the
        run spec and the run-directory name, so temperature-varying runs could not
        be located. Fix: locate run folders more robustly by reading the original
        run spec out of the generated artifacts rather than trusting the directory
        name.
\end{enumerate}

These three fixes are small but foundational: they are the first instances of a
theme that dominates the rest of the internship --- \textsc{helm}'s run identity is
carried by a drift-prone \emph{string}, and robust reproduction requires matching
on \emph{canonicalized recipe content}, not names (revisited at scale in
\S\ref{sec:canonkey} and in gotchas G1--G2, G10--G12; see
Appendix~\ref{app:gotchas}).

\section{Outcome and next steps recorded at the time}

The e2e tests established that HF and vLLM at \code{temperature=0} agree on the
core metric while the \code{temperature=1} control diverges as intended --- i.e.\
the instrument both confirms equivalence and detects real differences. The deck's
next steps (automate the e2e suite into CI, improve the warnings the systems emit,
and begin auditing more public models) foreshadow the rest of the summer.

%==============================================================================
\chapter{Corpus Cartography: Mapping the Reproducible Subset}
\label{chap:cartography}
%==============================================================================

\emph{Status deck of 2026-06-09.}

\section{The question: what \emph{can} be run at all?}

Before reproducing anything, one must know the shape of the corpus and which
slice is even a candidate. A statistics-and-filtering script was written to crawl
a directory of \textsc{helm} runs, extract per-run metadata, and emit a machine-
readable filter inventory (\file{filter\_inventory.json}). For each run it records
the model (parameter count, access class \{open, limited, closed\}), the scenario
and its class path, the benchmark family, whether the task requires a closed
(proprietary) judge, and a \emph{selection status} with explicit
\emph{failure reasons}. A representative row (a Thai-exam Typhoon-8B run) carries
fields such as \code{model\_num\_parameters}, \code{model\_access},
\code{model\_has\_hf\_client}, \code{size\_threshold\_params}, \code{eligible\_model},
\code{candidate\_pool}, and a \code{failure\_reason\_summary} --- the vocabulary
that the Stage-1 Sankey diagrams later consume.

\section{The numbers}

The crawl over the public corpus produced a stark funnel
(Table~\ref{tab:corpusstats}). Of \textbf{34{,}512} total runs, only
\textbf{1{,}109} were \emph{selected} as candidates for local reproduction under
the open-weight, runnable-recipe policy. The model-access distribution alone
explains much of the attrition: nearly as many \emph{limited}-access (18{,}273)
as \emph{open} (12{,}823) runs, plus 508 closed and 2{,}894 with no recorded
access class. Parameter counts were available for only 13{,}979 runs, so a size
budget could be applied to fewer than half the corpus with certainty.

\begin{table}[H]
\centering
\begin{tabular}{lr}
\toprule
\textbf{Quantity} & \textbf{Count} \\
\midrule
Total runs discovered            & 34{,}512 \\
Structurally incomplete          & 14 \\
Eligible model, out of scope     & 5 \\
\textbf{Selected (candidates)}   & \textbf{1{,}109} \\
\midrule
Model access: open               & 12{,}823 \\
Model access: limited            & 18{,}273 \\
Model access: closed             & 508 \\
Model access: none recorded      & 2{,}894 \\
\midrule
Runs with parameter-count data   & 13{,}979 \\
\midrule
Full grid points (models $\times$ scenarios) & 78{,}690 \\
Full grid: models / scenarios    & 366 / 215 \\
Filtered grid points             & 2{,}013 \\
Filtered grid: models / scenarios & 33 / 61 \\
\midrule
Runs requiring a closed judge    & 484 \\
Runs excluded \emph{solely} for closed judge & 25 \\
\bottomrule
\end{tabular}
\caption{The public-\textsc{helm} corpus funnel as measured on 2026-06-09. The
runnable subset is roughly 3\% of the corpus, and grid coverage is sparse: 198
of 2{,}013 filtered grid points are actually populated (``likely not all
scenario/model pairs make sense'').}
\label{tab:corpusstats}
\end{table}

\finding{The closed-judge dependency is real but not dominant.} Of 484 runs that
require a closed judge, only 25 are excluded \emph{solely} for that reason --- the
rest fail an independent filter too. This number motivated the later
\emph{open-judge extension} (\S\ref{sec:judge}): admitting closed-judge benchmarks
behind a flag and re-running them with an open judge recovers a
non-trivial-but-bounded slice.

\finding{Grid coverage is sparse and non-rectangular.} The full model$\times$scenario
grid has 78{,}690 points but only 6{,}844 are covered; the filtered grid has 2{,}013
points and 198 covered. This is the empirical basis for treating reproduction as a
\emph{coverage problem} --- the denominator of any ``fraction reproducible'' claim
must be the set of populated, in-scope grid points, not the Cartesian product. It
directly motivates the three-level coverage funnel (logical / recipe-canonical /
recipe-identical) that the pipeline reports (\S\ref{sec:funnel}).

\section{Why this is Stage 1}

This cartography is Stage 1 of the pipeline
(\file{index\_historic\_helm\_runs.py}): it discovers runs, applies the eligibility
filters, and emits the filter report whose Sankey diagram shows what was kept or
dropped and \emph{why}. Its output is the ``denominator'' the architecture
document (ADR-6) insists must remain explainable: every excluded run carries its
typed exclusion reason, so a reviewer can audit the claim ``we reproduced $X$ of
$Y$ runnable runs.'' Crucially, the exclusion reasons here are almost all
\emph{recipe/environment} reasons (\S\ref{sec:distinction}) --- gated model,
closed judge, no local deployment, over the size budget --- not reproducibility
verdicts.

\begin{quote}\small
\textbf{Corpus caveat (recorded later).} The curated candidate set that drives the
reproduction runbooks (\file{configs/run\_details.yaml}) is a hand-selected
Pythia/Vicuna/Qwen/OLMo subset, not the full eligible list; and an on-disk
\file{filter\_inventory.json} captured at one moment can go stale relative to the
live corpus (e.g.\ a modern-only pass that contains zero classic-era rows). The
distinction between ``the corpus'' and ``the curated candidate set'' matters when
reading any coverage number.
\end{quote}

%==============================================================================
\chapter{The Unified Comparison Core}
\label{chap:core}
%==============================================================================

\emph{Commits: 2026-06-11 through 2026-06-12 (roughly 30 commits); the ``Phase 0--3''
refactor and the ``Phase 3 / 4.0--4.9'' comparison-core unification.}

Two days of intense work rebuilt the analysis substrate so that every subsequent
reproducibility verdict flows through a single, framework-free comparison engine.
This is infrastructure, but it is infrastructure with research consequences: it is
what makes the later findings \emph{auditable} rather than asserted.

\section{The refactor (Phases 0--2): making the instrument legible}

The repository had grown to $\sim$86 modules and $\sim$32.8k lines with 18 console
scripts, and several ``god modules'' had accreted --- notably
\file{build\_reports\_summary.py} at 5{,}369 lines / 108 functions. The refactor
proceeded in phases: root hygiene and finishing an in-flight
\code{vllm\_service}$\rightarrow$\code{infer\_stack} rename (Phase 0); unifying the
CLI surface so every console script resolves to a thin \code{eval\_audit.cli.*}
wrapper (Phase 1); and decomposing the god modules into cohesive subpackages
(Phase 2).

The methodological point worth recording for the paper's ``reproducible software''
story is \emph{how} the splits were verified. Rather than hand-moving functions,
the god modules were split programmatically: AST-parse the module, compute the
top-level-symbol reference graph, validate the cluster assignment is a DAG,
generate each submodule with exactly the imports its symbols reference, then assert
that all top-level symbols are \emph{AST-identical} to \code{HEAD}. This gives a
machine-checkable ``pure relocation'' guarantee. A recurring subtlety --- that
monkeypatching a facade module no longer reaches functions that bind their
collaborators in a new home module --- caused several test fixes and is a lesson
about characterization tests surviving refactors.

\section{Phase 3 / \code{NormalizedDiff}: one core, many adapters}
\label{sec:normalizeddiff}

The deeper change was to unify two forked comparison cores. The \textsc{helm} path
computed agreement and diagnosis through \code{HelmRunDiff} (a 1{,}902-line class);
the \eeeshort path computed overlapping numbers through \code{normalized/compare.py}.
Both were invoked from one function, \code{\_build\_pair}, gated by a
\code{skip\_diagnosis} flag. The agreement math was already framework-free on both
paths; only the \emph{diagnosis} block and the \file{run\_spec.json} semantic diff
were \textsc{helm}-specific.

The unification introduced \code{eval\_audit/normalized/diff.py::NormalizedDiff}:
one framework-free core consuming two \code{NormalizedRun} objects and producing
agreement rows, agreement curves, quantiles, the facts-grade diagnosis, and the
judge-dependence metric-class split. The \textsc{helm} and \eeeshort paths became
\emph{thin adapters} over this one core; \code{HelmRunDiff}'s \code{\_diagnose\_repro}
was reduced to delegating to \code{normalized.diagnose.diagnose\_repro}.

The design was executed \emph{additive-first}, which is the load-bearing rigor:
\code{NormalizedDiff} was built next to the incumbent and proven \emph{byte-equal}
to committed baseline snapshots (fixtures F1--F10, gated at numerical tolerance
\code{atol}$=10^{-9}$, with the explicit rule that a looser tolerance is ``a
finding to investigate, not a knob to widen'') \emph{before} any renderer was
flipped to use it. Inverting the usual order --- build the replacement, prove
equivalence, then swap --- turned a ``high-risk hinge'' into two routine commits.

\section{The framework-free taxonomy and recipe-facts resolver}
\label{sec:taxonomy}

Three supporting modules were lifted out of the \textsc{helm}-specific package so
they carry no framework dependency:

\begin{itemize}[nosep]
  \item \file{eval\_audit/metrics\_taxonomy.py} --- classifies every metric as
        \emph{core} vs \emph{bookkeeping} (the metric-scale caveat from
        \S\ref{sec:inherited}) and, new here, as \emph{judge-dependent} vs
        \emph{deterministic} (\code{classify\_judge\_dependence}). The latter
        split lets a report separate the part of a comparison that is a
        \emph{reproduction control} (deterministic metrics) from the part that is
        a \emph{substitution measurement} (judge-dependent metrics).
  \item \file{eval\_audit/normalized/recipe\_facts.py} --- one resolver answering
        ``what recipe produced this run?'' in priority order: a native
        \code{recipe\_facts} block (JSON-encoded inside an \eeeshort aggregate's
        \code{source\_metadata.additional\_details}) $\rightarrow$ a sidecar
        \file{run\_spec.json} $\rightarrow$ \emph{unknown} (facts collapse to
        \code{status='unknown'} and the pipeline emits
        \code{comparability\_unknown:*} warnings). This is the seam through which
        the later ``unrecorded parameters'' story (Appendix~\ref{app:unrecorded})
        could be closed without a schema change.
  \item \file{eval\_audit/normalized/diagnose.py} --- the single input-to-label
        diagnosis implementation, including declared-substitution re-labeling
        (\code{intended\_substitution:*}).
\end{itemize}

\section{The open-judge extension (R2)}
\label{sec:judge}

A one-hour ``spike'' against the real corpus reshaped a planned feature and
recorded a genuinely useful finding
(\file{docs/planning/judge-identity-inventory.md}). The question was: where does
judge-model identity live for the six closed-judge benchmarks? Three findings:

\begin{enumerate}[nosep]
  \item \textbf{Judge identity is not in \file{run\_spec.json}.} Annotators carry
        \emph{empty} args; the judge model is hard-coded inside the \textsc{helm}
        annotator class as a function of \textsc{helm} version (e.g.\ a GPT-4o +
        Llama-3.1-405B ensemble). Officials therefore need a curated
        (annotator class, version)$\rightarrow$judge-models map, implemented as
        \file{eval\_audit/judge\_registry.py}.
  \item \textbf{The official ensemble already contains an open judge.} Because
        \textsc{helm} records per-judge sub-scores (e.g.\ \code{safety\_gpt\_score},
        \code{safety\_llama\_score}) as separate metrics, a \emph{same-judge
        reproduction control already exists inside the official data} --- one can
        compare \code{safety\_llama\_score} official-vs-local without any new run.
        This reframes substitution as per-metric rather than per-run.
  \item \textbf{Honesty scoping.} The \code{same\_judge} comparability fact is
        emitted \emph{only} for declared-substitution comparisons; emitting it
        everywhere would spray \code{comparability\_unknown:same\_judge} noise over
        every legacy pair. A crucial design decision recorded here: facts stay
        honest (\code{same\_judge: no} is never re-labeled to a ``substituted''
        status); only the \emph{diagnosis} re-labels the difference as
        \code{intended\_substitution:judge}. A ``substituted'' fact status was
        rejected as dishonest-by-construction.
\end{enumerate}

Stage 1 gained \code{--allow-closed-judge-benchmarks}, admitting these runs through
a distinct \code{judge-substitution} candidate pool that rides
\code{run\_details}$\rightarrow$manifests$\rightarrow$audit index to the planner.

\section{Why the entry points stayed separate}

A tempting simplification --- collapse the \textsc{helm} and \eeeshort CLIs into
one auto-detecting command --- was deliberately \emph{rejected}, and the reasoning
is worth preserving because it recurs. The \eeeshort-only path must remain
physically capable of building its full report tree \emph{with zero \textsc{helm}
artifacts on disk} (the ``EEE-only operability'' gate), because one of the paper's
case-study claims is that the \eeeshort per-instance schema alone suffices for
reproducibility analysis --- a claim a reviewer will want to verify by
\code{grep}, not by trusting a code review. Merging the CLIs would import the
\textsc{helm} adapter into the \eeeshort path and destroy that grep-checkable
guarantee.

%==============================================================================
\chapter{The OLMo Campaign}
\label{chap:olmo}
%==============================================================================

\emph{Commits: 2026-06-12 onward; status decks of 06-16, 06-23, 06-30, 07-08,
07-14. The OLMo reproduction is the internship's central case study, and the
vehicle through which every piece of reproduction machinery was hardened.}

\section{The target grid}

Six open-weight OLMo-family models were chosen as the reproduction subjects
(Table~\ref{tab:olmogrid}). Four are recent instruct models evaluated on the
capabilities benchmarks (\code{bbq}, \code{gpqa}, \code{ifeval}, \code{mmlu\_pro});
two are older base models with broad classic-track coverage. The grid spans both
\emph{no-chain-of-thought / temperature=0} benchmarks (bbq, commonsense, gsm,
ifeval, legalbench, med\_qa, mmlu, narrative\_qa, wmt\_14) and \emph{chain-of-thought
/ temperature=1} benchmarks (gpqa, mmlu\_pro).

\begin{table}[H]
\centering
\small
\begin{tabularx}{\textwidth}{lY}
\toprule
\textbf{Model} & \textbf{Benchmarks reproduced} \\
\midrule
\code{olmoe-1b-7b-0125-instruct}   & bbq, gpqa, ifeval, mmlu\_pro \\
\code{olmo-2-1124-7b-instruct}     & bbq, gpqa, ifeval, mmlu\_pro \\
\code{olmo-2-1124-13b-instruct}    & bbq, gpqa, ifeval, mmlu\_pro \\
\code{olmo-2-0325-32b-instruct}    & bbq, gpqa, ifeval, mmlu\_pro \\
\code{olmo-7b}                     & mmlu (62), math (7), legalbench (5), wmt\_14 (5), \\
                                   & natural\_qa (2), commonsense (1), gsm (1), \\
                                   & med\_qa (1), narrative\_qa (1) \\
\code{olmo-1.7-7b}                 & mmlu (57) \\
\bottomrule
\end{tabularx}
\caption{The six OLMo reproduction subjects and their benchmark coverage. The
instruct models exercise the chat-template and judge-dependent paths; the base
models exercise the classic-track, completions path.}
\label{tab:olmogrid}
\end{table}

Early smoke tests already surfaced two environment issues that recur:
a shared-machine configuration-file conflict, and \code{olmo-7b}'s tokenizer
attempting to run external code (an early symptom of the EOS-append bug diagnosed
in \S\ref{sec:eosappend}). These were fixed via HELM tokenizer-alias corrections,
adding \code{sentencepiece}/\code{tiktoken} tokenizer backends as dependencies, and
loading the OLMo-7B tokenizer through a HELM entry-point plugin.

\section{Chat-template overhead blows the context window}

The first substantive reproducibility bug in the OLMo runs was a
\code{ContextWindowExceededError}: the chain-of-thought runs
(\code{num\_output\_tokens=2048}) on the 4096-context chat models were rejected by
vLLM. The mechanism is subtle and paper-worthy. \textsc{helm}'s
\code{LocalWindowService} truncates the prompt to
$\text{max\_seq\_len} - \text{expected\_completion} = 4096 - 2048 = 2048$ tokens,
but it counts the \emph{raw} prompt, \emph{before} the chat template is applied.
When the run is routed through vLLM's OpenAI-compatible chat endpoint, the OLMo-2 /
OLMoE chat template wraps the prompt and adds tokens \textsc{helm} never counted.
Measured with the actual tokenizers, the single-turn wrapper overhead is
\textbf{12 tokens (OLMo-2)} and \textbf{13 tokens (OLMoE)}, so
$2048 + 13 + 2048 = 4109 > 4096$ and vLLM hard-rejects.

The fix reserved a 32-token margin per preset via the sanctioned
\code{helm\_max\_sequence\_and\_generated\_tokens\_length} knob (already used by the
Vicuna chat path), and repaired a latent bug where the flat-preset code path
silently dropped that override. The general lesson --- \emph{when \textsc{helm}
drives a chat model through an OpenAI-compatible server, its prompt-token budget is
blind to the chat template the server applies} --- is the first of several
chat-template-versus-window interactions, and it recurs when the Qwen turbo
endpoints hit the same wall in July (\S\ref{sec:qwenwindow}).

\section{Data-access failures are filtering reasons, not results}

Several OLMo benchmarks could not be run at all, and classifying them correctly was
the point:

\begin{itemize}[nosep]
  \item \code{competition\_math} was disabled on the HF Hub (script-based dataset,
        401 on the loader).
  \item \code{natural\_qa}: the public GCS bucket revoked \emph{authenticated}
        reads (it dropped \code{allAuthenticatedUsers}, not just \code{allUsers}),
        so no credential unblocks it. An entire staging shim built to work around
        the assumed ``authenticated reads work'' had to be reverted once the
        access assumption was falsified by a single test with real credentials.
        The recorded lesson --- \emph{validate the access assumption before
        building the plumbing that depends on it} --- is a methodological point for
        the paper's failure taxonomy.
  \item \code{gpqa} is a gated HF dataset (requires accepting terms).
\end{itemize}

Each was disabled as a documented \emph{environment} failure, joining a taxonomy
that is deliberately kept separate from the reproducibility verdicts.

\section{The canonical-key matching bug: 114 phantom misses}
\label{sec:canonkey}

The most consequential ``bug'' of the OLMo campaign was not in execution but in
\emph{pairing}. The planner matched logical-key \emph{strings} by generating
separator permutations and (under a flag) stripping \code{groups=} tokens, but it
never canonicalized \emph{token order}. The OLMo MMLU keys are the same token set
in a different order, so the local and official variant sets had an \emph{empty
intersection}: all 114 MMLU packets were reported as
\code{missing\_official\_component} even though the public counterparts existed in
the index.

The fix (\file{docs/planning/core-report-planner-robust-matching-plan.md}) replaced
the order-sensitive matcher with a single \code{canonical\_logical\_key} in
\file{eval\_audit/run\_entries.py}: parse the \code{benchmark:k=v,...} name, drop
bookkeeping-only tokens (\code{groups}, \code{model\_deployment}), canonicalize
values (model \code{/}$\leftrightarrow$\code{\_}, \code{mmlu\_pro} \code{subject}
$\rightarrow$ \code{subset}), and re-serialize with the key--value pairs
\emph{sorted by key}. Once you sort, the separator-permutation and groups-strip
variants all fold into one deterministic string, so keeping both matchers would
have re-created the very ``two competing normalizers'' divergence the fix set out
to end. The real-data validation expectation: \code{n\_skipped} $114\rightarrow\sim0$,
\code{n\_built} $35\rightarrow\sim149$.

This is a symmetric equivalence, correct for \emph{grouping}; the asymmetric subset
test (\code{run\_dir\_matches\_requested}) was deliberately preserved for
\code{compare\_batch}, which answers the different question ``does this candidate
satisfy this request.'' The distinction --- symmetric equivalence for grouping vs
asymmetric subset for satisfaction --- is a reusable design point.

\section{The BBQ prompt drift: a genuine recipe disagreement}

Not every difference is an artifact to be fixed. The BBQ runs surfaced a
\emph{real} recipe drift: the official run entry carried
\code{output\_format\_instructions=mcqa} and the prompt began ``Answer with only a
single letter.'' (followed by the standard preamble) while the locally
reconstructed run entry dropped that token and produced only ``The following are
multiple choice questions (with answers).'' The model sees a different prompt and
produces measurably different output. This is gotcha G5 (\S\ref{app:gotchas}), and
it is precisely the kind of difference that must be \emph{surfaced} (via the
\code{same\_instructions} comparability fact), not silently fixed --- and it is one
of the strongest motivations for the from-spec replay of
Chapter~\ref{chap:fromspec}: replaying the official \file{run\_spec.json} verbatim
makes such reconstruction drift vanish by construction.

%==============================================================================
\chapter{Containerization and GPU Leasing}
\label{chap:containers}
%==============================================================================

\emph{Commits: 2026-06-16 through 2026-06-24. Two intertwined infrastructure
efforts: pinning the execution environment in Docker, and letting many HELM runs
fan out across GPUs under a lease.}

\section{Why containerize: the environment is a variable}

By default \code{eval-audit-run} executes each \textsc{helm} run-entry in the host
venv of whatever GPU machine the scheduler lands on, making torch / CUDA /
transformers / \textsc{helm} versions an \emph{uncontrolled variable} --- exactly
the kind of drift the audit exists to separate from true reproducibility failures.
The containerized path runs each run-entry inside a pinned Docker image and records
\emph{which image, by \code{sha256} digest, produced each run}.

The design (confirmed interactively) is worth recording:

\begin{itemize}[nosep]
  \item A brand-new, independent multi-stage Dockerfile (CUDA devel$\rightarrow$
        runtime, \code{uv}-managed standalone CPython, venv at \code{/opt/venv}),
        built from \emph{pristine committed submodule state} via \code{git archive}
        at each submodule's gitlink SHA --- no remote dependency, no \code{.git}
        leakage, SHAs recorded as OCI labels.
  \item The image tag is resolved to an immutable \code{<repo>@sha256:<digest>}
        reference \emph{once at schedule time} and every node is pinned to it, so a
        rebuild automatically re-pins.
  \item The container runs as root; a thin entrypoint chowns \emph{only} the output
        dir back to \code{HOST\_UID:HOST\_GID} on exit, so \code{kwdagger} on the
        host can own results while in-container HF-cache writes still succeed.
  \item The docker-wrapping \code{kwdagger} node \emph{subclasses} the magnet
        materialize node and overrides only \code{command}, inheriting the
        \code{out\_paths}/\code{DONE}-sentinel contract for free --- the reusable
        lesson being ``wrap an existing scheduler node in a new execution substrate
        by subclass + override, not reimplementation.''
\end{itemize}

Several first-real-build bugs are instructive because each could only surface at
the first genuine build or run and each maps to a general lesson:

\begin{itemize}[nosep]
  \item \textbf{Dangling interpreter symlink} --- \code{uv}'s managed standalone
        Python was not copied into the final stage, so \code{/opt/venv/bin/python}
        was a dangling symlink reporting ``command not found.'' Lesson:
        \emph{uv-managed Python + multi-stage Docker = copy the interpreter, not
        just the venv}, and put the import smoke-test in the stage that ships.
  \item \textbf{HF token not reaching the container} --- \code{kwdagger} runs each
        job in a fresh tmux pane whose environment is deliberately empty (secrets
        hygiene), so \code{-e HF\_TOKEN} forwarded nothing; the host-venv path had
        silently relied on the on-disk \code{huggingface-cli login} token. Fix:
        write the resolved token to \code{<hf\_cache\_dir>/token} (mode 0600) at
        schedule time, in the process that \emph{did} inherit the env. Lesson:
        \emph{a redundant channel that silently carries the load hides the breakage
        of the primary one}.
  \item \textbf{Wrong extras and an unpinned hub} --- the image installed
        \code{crfm-helm[heim]} (image-generation extra) which lacks
        \code{langdetect} (needed by ifeval), and left \code{huggingface\_hub}
        unpinned so it floated to a version whose repo-resolution API broke
        \code{wmt\_14}. Fix: install \code{crfm-helm[all]} and co-pin
        \code{huggingface\_hub==0.36.2}, asserted in \emph{both} the install layer
        and the final stage. Lesson: \emph{pin the secrets-of-resolution
        (transitive floats), not just the top-level package}, and pair a build-time
        guard with a runtime probe of the resolved artifact (a stale pinned digest
        is invisible to the build guard).
\end{itemize}

\section{From profiles to a leasing catalog}

Mid-June the \code{infer\_stack} submodule was rewritten from a
profile/contract model to a \emph{catalog} (models + endpoints) plus \emph{leasing}
(\code{acquire}/\code{release}) model. The two in-scope runbooks (phi-2 e2e and the
six OLMo presets) were migrated: \code{load\_profile\_contract} was replaced with a
pure static \code{resolve\_serving\_facts} over the catalog; the config dirs were
reschema'd from \code{config.yaml}+\code{models.yaml} to
\code{settings.yaml}+\code{catalog.yaml}; and an explicit \code{protocol\_mode}
(completions for base models, chat for instruct) was made a required preset fact
rather than a silent default. A subsequent question --- ``do \code{serve} and
\code{acquire} both need to exist?'' --- collapsed the two verbs into
\code{acquire} once a one-line grep proved the distinguishing \code{owner='manual'}
field drove no behavior.

\section{The leasing pipeline: preconditions as job phases}

The deeper design question was how to drive \emph{high-parallelism} reproduction:
each \textsc{helm} run is a \code{kwdagger} \code{ProcessNode} that must bracket
itself with a GPU lease (acquire before, release after), where release must run
after finish \emph{or} crash. The decisive argument, reached across three rounds of
design discussion, was \emph{co-location}: the lease lives in a node-local ledger
and the gateway is localhost, so release must run \emph{where} acquire ran. A
separate cleanup DAG node cannot promise that (and is actively wrong on multi-node
scheduling), and success-only dependencies (\code{afterok}) skip a release node
when the run fails --- the exact leak to avoid. The right primitive is a job
\emph{phase} (setup/teardown), which is also the more general, reusable scheduler
feature.

The findings that settled the design were verified against source, not assumed:
\code{cmd\_queue}'s \code{preamble} already gates the main command, so only an
always-run \code{teardown} trap was new; \code{kwdagger}'s cache guard
(\code{test -e DONE || cmd}) means a \emph{skipped} job acquires nothing; and a
blocking acquire's retry loop \emph{reclaims TTL-expired leaks while it waits}, so
queue-and-wait and leak-recovery become one mechanism --- \emph{iff every pipeline
lease has a finite TTL}. A subtle bash-precedence trap was caught and closed: a
setup chain \code{acquire \&\& \{ snapshot \} || true} lets a \emph{failed}
acquire fall through to \code{|| true}, defeating the gate; the fix scopes
\code{|| true} inside the snapshot brace so only the snapshot is best-effort.

Two further design moves are worth recording because they generalize:

\begin{itemize}[nosep]
  \item \textbf{Decoupling leasing from containerization.} The lease bracket was
        initially bolted onto the docker node only, which would have dragged
        containerization onto the host-venv vLLM scenarios. It was extracted into a
        transport-agnostic \code{LeaseBracketMixin} + shared command renderer;
        leasing and containerization became \emph{orthogonal} axes (a run can be
        docker+lease, docker+no-lease, or --- briefly --- bare+lease). The lesson:
        when a capability seems ``coupled'' to a transport, check whether it is the
        capability that needs the transport or just where someone bolted it on.
  \item \textbf{Scoped cleanup over global teardown.} Every \code{release --all
        --evict} was removed from the runnable path in favor of per-run release
        (\code{reclaim: stop}) plus a scoped \code{infer-stack gc} that reclaims
        only TTL-expired leaks in this user's ledger --- so the shared docker stack
        is never torn out from under a co-tenant. A bootstrap that exists only to
        read a managed secret pays for it with a scoped \code{--env-file} handle,
        not a global teardown.
\end{itemize}

By the end of this arc, containerization was made \emph{mandatory} and the
non-containerized code paths were deleted, with leasing the orthogonal opt-in ---
so the execution environment is always a pinned, digest-recorded artifact
(\code{container\_provenance.json}).

%==============================================================================
\chapter{Faithful Replay: From Run-Entry to Run-Spec}
\label{chap:fromspec}
%==============================================================================

\emph{Commits: 2026-06-25 through 2026-06-30; the \code{impl/run-from-run-spec}
branch on which most subsequent work sits.}

\section{The core methodological principle}

This chapter contains what is arguably the single most important methodological
decision of the internship, and it should anchor the paper's ``methods'' section.

Public \textsc{helm} runs do not record the CLI run-entry string that produced
them, nor the exact \textsc{helm} version. What they \emph{do} ship is a fully
resolved \file{run\_spec.json}: the complete adapter spec, scenario spec, metric
specs, and annotators, with every default already materialized. The prior approach
--- and the approach every earlier reproduction took --- was to \emph{reconstruct}
a run-entry string from the run and re-parse it. But re-parsing re-derives the
recipe under \emph{today's} library defaults, which silently drifts from the
recipe that actually produced the official numbers (this is the mechanism behind
gotchas G1, G5, and the BBQ drift of \S\ref{sec:canonkey}).

\begin{quote}\itshape
The principle: never reconstruct the recipe. Replay the resolved
\file{run\_spec.json} \emph{verbatim}.
\end{quote}

Concretely, a new \code{materialize\_helm\_run\_from\_spec.py} CLI (a faithful-
replay sibling of the run-entry materializer) reuses the discovery/matching
scaffolding but swaps the compute step from ``reconstruct a run-entry string
$\rightarrow$ \code{helm-run} subprocess'' to ``\code{from\_json(run\_spec.json)}
$\rightarrow$ in-process \code{run\_benchmarking}.'' In discovery mode, the
run-entry string now does nothing but \emph{locate} the official directory; the
authoritative resolved recipe drives execution. This memory --- \emph{all
reproductions must be from-spec} --- became a standing rule, and the older OLMo
run-entry path was flagged for migration.

\section{Two failure modes, one loud and one silent}

The design carefully distinguishes the two ways a from-spec replay can go wrong,
because they demand opposite handling (this is gotcha-adjacent and reviewer-
relevant):

\begin{description}[style=nextline,leftmargin=1.25em]
  \item[Class drift is loud.] If the \file{run\_spec.json} references a scenario or
        metric class that does not resolve under the local \textsc{helm}
        (\code{get\_class\_by\_name} raises \code{ImportError}), the preflight
        catches it before any instance runs. This is the correct behavior: a class
        that cannot be resolved is an environment mismatch, and it is surfaced as
        such. (The preflight's exception net had to be broadened from
        \code{(ImportError, AttributeError)} to \code{Exception}, because importing
        a vision--language scenario raises \textsc{helm}'s
        \code{OptionalDependencyNotInstalled} --- itself a recipe/environment
        filtering reason, not a crash to propagate.)
  \item[Field drift is silent.] The cattrs codec ignores unknown keys and fills
        defaults, so a dropped field produces no error --- only wrong numbers.
        The guard is a round-trip test: because \textsc{helm} \emph{writes}
        \file{run\_spec.json} with \code{asdict\_without\_nones} (dropping only
        \code{None} values), every on-disk key is non-\code{None}, and the
        re-serializer re-emits every non-\code{None} field --- so a key missing
        after \code{from\_json}$\rightarrow$\code{to\_json} is an unambiguous silent
        drop. Ten sampled MMLU specs round-tripped with zero key loss.
\end{description}

The reusable insight recorded at the time: \emph{writer/serializer asymmetry
decides whether a diff test is sound} --- the no-key-dropped assertion holds only
because the writer drops a superset of what the re-serializer drops.

\section{Making the deployment substitution visible}

A from-spec replay serves the model on local hardware, which is by definition a
\emph{different deployment} than the original (e.g.\ a hosted Together API). The
subtle danger is that a pure by-name replay would record the \emph{official}
deployment name in the produced run, so the comparison would report
\code{same\_deployment=yes} and the single most important difference the audit
exists to surface --- that the serving engine changed --- would become invisible.

The fix (\file{from-spec-deployment-rewrite-plan.md}) rewrites
\code{adapter\_spec.model\_deployment} to the \emph{local} name after
deserialization, via an explicit \code{--model-deployment} argument. Explicit,
not auto-derived: the whole feature exists to stop \emph{silently} masking the
substitution, so its own failure mode must be loud --- a wrong name is a
\textsc{helm} ``deployment not found'' crash \emph{before any instances run},
whereas auto-derive's failure mode is the silent no-op that fails \emph{toward} the
very bug it fixes. The produced run then records the served endpoint, and the
existing diff logic correctly reports \code{same\_deployment=no} with no downstream
plumbing.

\section{Addressing runs by (root, relative-path) and content-addressed caching}

A further refinement replaced the drift-prone run \emph{name} as the addressing
mechanism with \emph{host-side} resolution by \code{(precomputed\_root,
rel\_path)}: resolve the official \file{run\_spec.json} at schedule time, apply
\emph{raw-JSON} scalar edits (only \code{adapter\_spec.model\_deployment} and
\code{max\_eval\_instances}, so no cattrs round-trip and therefore no field drift),
and materialize a substituted copy under a staging directory. Each run rides one
\code{kwdagger} \emph{submatrix} entry (the verified zip primitive) so the per-run
tuple travels together without an $N\times N$ cross-product. Two consequences worth
recording:

\begin{itemize}[nosep]
  \item \textbf{No image rebuild, no corpus mount.} The runner needs only the tiny
        staging directory, because substitution happens host-side.
  \item \textbf{Content-addressed caching.} \code{max\_eval\_instances} is baked
        into the materialized copy but is not a matrix parameter, so a naive fixed
        filename (\code{run\_spec.json}) would make \code{kwdagger} \emph{skip} and
        reuse a stale run after a cap change. The fix appends a content hash to the
        filename: \code{run\_spec.<sha256(bytes)[:16]>.json}. Because the
        materialized path \emph{is} the node's algorithm identity, an identical
        recipe yields the same path (skip) and any changed field yields a new path
        (recompute). This is the reusable ``content-address generated input files
        so baked-in changes invalidate the cache'' rule.
\end{itemize}

An exporter ``auto-freeze'' resolves each preset run-entry to its exact rel-path
\emph{once} at export time (against the corpus snapshot) and freezes the sources
into the generated manifests --- the only remaining place token-subset discovery
runs, pinned where drift cannot occur --- and a dry-check existence mode validates
that a frozen manifest's rel-paths still resolve.

%==============================================================================
\chapter{Deployment Matching: The Flagship Investigation}
\label{chap:deploymatch}
%==============================================================================

\emph{Commits: 2026-06-30 through 2026-07-10; the \code{dev/tools/deployment\_match/}
tool and the documents \file{vllm-vs-huggingface-deployment-match.md} and
\file{helm-unrecorded-deployment-params.md}. This is the scientific heart of the
internship: the sustained effort to attribute residual disagreement to a specific,
unrecorded execution parameter.}

\section{The provocation: OLMo-7B produces prompt-independent gibberish}
\label{sec:eosappend}

The investigation began with a dramatic symptom. During from-spec reproduction of
\code{narrative\_qa}, local OLMo-7B repeated the same prompt-independent
boilerplate --- ``The first thing you need to do is to make sure that you have
a\ldots before you start playing.'' --- for \emph{every} instance, while the
official \code{together/olmo-7b} conditioned on the prompt correctly. The diff
showed \code{prompts\_equal=True} (the prompt reached the model) and a hugely
negative local completion log-prob against the official's $0$: the model ran but
ignored its input.

The first hypothesis was a documented one: OLMo-1 7B \code{float16} numerical
instability (activation outliers exceed fp16's range $\rightarrow$ inf/NaN in
attention $\rightarrow$ collapse to the highest-prior pretraining attractor). A
minimal deployment matrix was built to test it --- and the hypothesis was
\emph{refuted}. The true cause was a \emph{tokenizer} issue: \code{OLMo-7B-hf}'s
tokenizer appends an \code{<|endoftext|>} token to the prompt under vLLM's default
\code{add\_special\_tokens=True}; the Together deployment did not do this, and
\textsc{helm} never recorded the difference. The production fix serves the
OLMo-1.7 tokenizer (whose post-processor omits the append), and the completions
path forwards \code{add\_special\_tokens=false}.

\finding{This is the paradigm case of the whole project.} A disagreement that
looked like flat irreproducibility (the model outputs nonsense) was in fact a
\emph{recipe/environment} failure caused by a single unrecorded request-time
tokenizer flag. And the first, plausible theory (fp16 instability) was wrong ---
which is exactly why the response was to build a \emph{sweep} rather than trust a
single-knob theory. The lesson --- that single-knob theories about deployment
divergence had by now been wrong twice --- became the design charter for the
deployment-match tool.

\section{The tool: sweep, don't guess}

\code{dev/tools/deployment\_match/} takes \emph{any} public \textsc{helm} run,
builds a grid of plausible local deployments for that run's model, evaluates a
small length-stratified prompt sample, and ranks each cell by closeness to the
official completions. Because Together stores no log-probs, the objective is
\emph{agreement with the official completion text} (a composite of quasi-exact-
match, first-word match, and similarity, with a model-agnostic
``prompt-independence collapse'' detector). The grid is two-tier: serve-time knobs
(dtype $\times$ tokenizer $\times$ \code{max\_model\_len}) become distinct
infer-stack endpoints; request-time knobs (\code{add\_special\_tokens},
\code{add\_generation\_prompt}, protocol) are probed per request against one
container. The tool has \code{dry-run} / \code{run} / \code{confirm} / \code{auto}
phases, and \code{confirm} produces a real one-endpoint catalog you can
\code{acquire}.

A key design principle, formalized after the fp16 mis-diagnosis, splits ``exclude
settings beforehand'' into two opposite treatments:
\begin{itemize}[nosep]
  \item \textbf{Feasibility pruning (safe).} A recipe that physically cannot serve
        on this hardware teaches nothing and wastes a full serve cycle, so it is
        dropped a priori with a typed reason (mirroring the Stage-1 filter
        pattern). Example: \code{float32} on a MoE model $\rightarrow$
        \code{infeasible:moe-fp32-shared-mem} (see \S\ref{sec:fp32}).
  \item \textbf{Relevance pruning (risky) is refused.} ``This knob probably won't
        change the output'' is exactly what the tool will not guess; such knobs
        stay in the sweep unless the operator narrows an axis explicitly.
\end{itemize}

\section{Chat-template version drift: \code{add\_generation\_prompt}}
\label{sec:addgenprompt}

The next finding is a subtle, real reproducibility trap on OLMoE-instruct, and it
is a clean example of a \emph{software-version} dependency hiding inside a recipe.
\textsc{helm}'s \code{get\_prompt} always calls
\code{apply\_chat\_template(\ldots, add\_generation\_prompt=True)} --- but that flag
only does anything if the chat template \emph{implements} it
(\code{\{\% if add\_generation\_prompt \%\}<|assistant|>\textbackslash n\{\% endif \%\}}).
The OLMoE template shipped with \textsc{helm}'s \emph{older} transformers
\emph{ignored} the flag, so \code{add\_generation\_prompt} was effectively
\emph{False} and the model was fed the prompt \emph{without} the trailing
\code{<|assistant|>\textbackslash n}. Modern transformers/vLLM ship an updated
template that honors it, so the same call now \emph{appends} the suffix
$\rightarrow$ a different prompt $\rightarrow$ a different output on every cell.
Rendering with \code{add\_generation\_prompt=false} on modern transformers moves
the output back toward the official, confirmed empirically. The tool now sweeps
\code{add\_generation\_prompt} $\in \{$True, False$\}$ as a cheap request-time knob
and lets the scorer decide.

\section{The float32 discovery: the flagship result}
\label{sec:fp32}

The single most consequential finding of the internship concerns model
\emph{precision}, and it is the one most worth foregrounding in the paper.

\textsc{helm}'s \code{HuggingFaceClient} loads a model with only the kwargs the
deployment config supplies. For \code{huggingface/olmoe-1b-7b-0125-instruct}, the
config gives exactly \code{args: \{device\_map: auto\}} --- \emph{no}
\code{torch\_dtype}. So \code{AutoModelForCausalLM.from\_pretrained(\ldots,
device\_map="auto")} is called with no dtype, and \emph{every transformers 4.x}
defaults that to \code{torch.float32} for backward compatibility, \emph{ignoring
the checkpoint's bf16 config}. (Reading the checkpoint's dtype is a transformers
v5 change; even late 4.x still hits ``\code{else: set fp32 as the default dtype for
BC}.'') The OLMoE architecture landed in transformers $\sim$4.45, so a January-2025
run is squarely in the fp32-default regime. This is a \emph{deductive} claim ---
read off the loader source and the library's version-dependent default, not fit to
outputs --- so it is stronger than an observational match. Its residual is the one
parameter it depends on: the exact \code{transformers} version, which the run dir
does not record (only bounds). Under that reading:

\begin{quote}\itshape
The code path \textsc{helm} used loads the weights in \code{float32} under every
\code{transformers} version consistent with the OLMoE run's date --- so the
official most likely ran \code{float32}. (Conditional on \code{transformers}$<$5;
a run produced under \,$\geq$5 would default to bf16 and move this conclusion.)
\end{quote}

A small oracle-scored \emph{probe} corroborates this: reproducing the official on
the HF side at \code{float32} (\code{transformers.generate()}) reaches quasi/exact
match $\approx 1.0$ on the sampled instances, versus $\approx 0.17$ for the best
fp16 vLLM cell. Table~\ref{tab:fp32} shows the overnight probe sweeps across the
four OLMo instruct models. \textbf{This is deployment-\emph{matching} evidence, not
a completed HELM reproduction} (see the epistemic-status note below): the scores are
composites over $n\approx12$ ifeval instances, scored against the sampled official
completions, and the winning cells additionally rely on a probe-only
\code{add\_special\_tokens=False} request knob that an ordinary \textsc{helm} run
does not send.

\begin{table}[H]
\centering
\small
\begin{tabular}{lll}
\toprule
\textbf{Model} & \textbf{HF best} & \textbf{vLLM best} \\
\midrule
OLMoE-1B-7B (MoE)  & \textbf{fp32 MATCH 0.971} (quasi 1.0) & fp16 PARTIAL 0.494 \\
OLMo-2-7B (dense)  & fp32 PARTIAL 0.175 (quasi 0.0) & \textbf{fp32 MATCH 0.915} \\
OLMo-2-13B (dense) & fp32 PARTIAL 0.324 (quasi 0.0) & \textbf{fp32 MATCH 0.904} \\
OLMo-2-32B (dense) & fp16 PARTIAL 0.234 (quasi 0.0) & \textbf{fp32-tp2 MATCH 0.961} \\
\bottomrule
\end{tabular}
\caption{Overnight deployment-match sweeps ($n=12$ ifeval instances, scored against
the official completions). fp32 is the matching precision; the winning
\emph{engine} depends on architecture. OLMoE is MoE, so vLLM's Triton fused-MoE
kernel cannot serve fp32 (shared-memory limit) and only HF fp32 reproduces it; the
dense OLMo-2 models reproduce near-exactly under vLLM fp32.}
\label{tab:fp32}
\end{table}

\par\medskip\noindent\fbox{\parbox{0.97\textwidth}{\small
\textbf{Epistemic status of the fp32 result (interpretive revision, 2026-07-15).}
This table is a \emph{configuration-discovery} result, not a reproduction result.
Read it as: \emph{among the deployment configurations tested, \code{float32}
candidates produced the highest agreement with the sampled public OLMo outputs.}
It does \emph{not} yet establish ``the historical deployment was \code{float32}''
as an identified fact, nor ``\code{float32} reproduced the official \textsc{helm}
result,'' for three reasons: (i) matching a small oracle sample does not uniquely
identify the producing deployment --- multiple latent configurations can be
observationally equivalent on $n\approx12$; (ii) the winning cells use a probe-only
\code{add\_special\_tokens=False} knob not sent by an ordinary \textsc{helm} run;
(iii) the ``confirm'' step --- land the winning config on the real \textsc{helm}
execution path, run the exact public \file{run\_spec.json}, and compare through the
standard pair-report pipeline on \emph{held-out} instances --- has not been run for
any model. The \emph{deductive} dtype-default argument above is the stronger leg and
survives these caveats (it reads the loader, not the outputs), but it too is
conditional on the unrecorded \code{transformers} version. See
Appendix~\ref{app:revision} for the consolidated epistemic status and the corrected
disagreement taxonomy.}}

\par\medskip
Two aspects make this finding sharp:

\begin{itemize}[nosep]
  \item \textbf{The matching precision is the one cell that cannot run.} vLLM's
        Triton fused-MoE kernel needs $\sim$131~KiB of shared memory per block in
        fp32 --- over the $\sim$99~KiB cap on workstation cards --- so it OOMs at
        the \emph{kernel} level (not VRAM). Every runnable local vLLM cell for
        OLMoE is reduced-precision, and none reproduces an fp32 reference. That
        \emph{fp16 beats bf16} among those cells is itself corroboration: fp16's
        finer mantissa flips fewer greedy near-ties against a high-precision
        reference, which says the reference is higher-precision than bf16 ---
        consistent with fp32. Serving fp32 requires either an H100-class card (big
        shared memory) or, for dense models, tensor-parallel fp32
        (\code{--fp32-tensor-parallel-size 2}); for the MoE, HF \code{generate()}
        at fp32 is the only route.
  \item \textbf{It reframes the disagreement.} A bf16-vs-official OLMo-2
        disagreement is a \emph{deployment-boundary artifact} (a precision the
        local engine could not easily provide), \emph{not} a reproducibility
        failure --- re-running at fp32 closes it. This is precisely the
        recipe/environment-vs-reproducibility distinction of \S\ref{sec:distinction},
        applied to the most consequential single case.
\end{itemize}

\section{The residual puzzle: the HF probe is broken for dense OLMo-2}

Rigor demanded following up an inconsistency. The HF probe reproduced OLMoE
\emph{exactly} but gave near-zero quasi-match for the three dense OLMo-2 models,
even though vLLM fp32 reproduced them. The forensic diagnosis pinned it precisely:
tokenizing the exact ifeval prompt showed byte-identical tokens across every path
(so it is \emph{not} the prompt); yet the HF probe's first token disagreed with the
official (first-token agreement 0.42) while vLLM fp32 matched character-for-
character (first-token 1.00). A first-token disagreement on a byte-identical prompt
is a clean litmus: at step 0 there is no accumulated history, so a different first
token can only mean a different \emph{forward-pass} logit vector. The divergence is
a fp32 forward-pass difference, not a prompt or precision-class difference ---
attributable to knobs \textsc{helm} never records: \code{attn\_implementation},
device placement (\code{device\_map=auto} shards a fits-on-one-GPU model across
GPUs, changing fp32 reduction order), and the decode path (\textsc{helm}'s
\code{do\_sample=True, temperature=1e-7} vs the probe's greedy). The probe was
extended to sweep these forward-pass axes. \textbf{Interpretive revision: this is a
finding, not merely a broken probe.} An earlier framing dismissed the OLMo-2 HF
divergence as ``not a science problem'' (a defect in our tool, since vLLM fp32 did
match). The sharper reading, which the paper should adopt: a \emph{current}
\code{transformers.generate()} at fp32 does \emph{not} reproduce an
\emph{apparently} HF-produced historical result \emph{even under byte-identical
prompt tokens}. That is precisely an \emph{unresolved systematic divergence under
prompt-token equivalence} --- direct evidence of execution-stack sensitivity, and a
core datum for the thesis, not an artifact to wave away. Its cause is not yet
isolated (candidates: model/tokenizer revision, \code{transformers}/\code{torch}
version, attention implementation, generation defaults, device placement, numerical
precision, CUDA behavior). That vLLM fp32 happens to match is a separate fact and
does not make the HF$\to$HF gap uninteresting. The reusable insight stands and is
sharper for it: \emph{``same engine family'' (transformers$\rightarrow$transformers)
does not imply reproducible} --- device sharding, attention-implementation drift, and
sample-vs-greedy each perturb fp32 logits enough to diverge greedy decoding.

\section{Naming the gap: the unrecorded execution substrate}
\label{sec:substrate}

The investigation crystallized into a taxonomy (\file{helm-unrecorded-deployment-
params.md}, Appendix~\ref{app:unrecorded}): \textsc{helm} serializes the
\emph{recipe} and the model's \emph{name} but almost nothing about the
\emph{execution substrate}. Empirically, of 148 HuggingFaceClient deployments in
\textsc{helm}'s \code{model\_deployments.yaml}, only 19 pin \code{torch\_dtype}
(all \code{bfloat16}) and only 7 pin a model \code{revision}. Ranked by effect on a
greedy output, the unrecorded parameters are: (Tier 1) weight revision, load
precision, quantization, software-stack versions; (Tier 2) attention implementation,
matmul/TF32 precision, hardware, device topology, batch composition; (Tier 3)
tokenizer + chat-template version, \code{add\_special\_tokens}/BOS handling,
\code{generation\_config.json} defaults, stop-sequence/detokenization semantics.
The one high-leverage lever is \emph{pinning the model revision (commit SHA)},
because the tokenizer, chat template, and generation config all live inside the
model repo --- so \code{model\_revision} + \code{torch\_dtype} +
\code{transformers\_version} together close the majority of the gap. Crucially,
these must live in \code{client\_spec.args} / \code{catalog.yaml} / a
\code{recipe\_facts} block, \emph{not} in the run-spec name or
\code{model\_deployment} identity string --- a SHA in the pairing key would break
official$\leftrightarrow$local pairing, whereas a SHA in the provenance block is
invisible to it.

A companion effort (\code{hf\_inprocess.py} + infer-stack
\code{acquire --reserve-gpus N}) built the mechanism to reproduce a
HuggingFaceClient official \emph{the same way it was produced} --- in-process
\code{transformers.generate()} at pinned fp32, on a GPU held (but not served) by a
reserve-only lease. The mechanism landed; the routing switch that would make the
default replay path choose it for a HuggingFaceClient official was scoped but left
unwired (a replay still defaults to vLLM).

%==============================================================================
\chapter{Era-Pinned Containers: Reproducing the Pre-v0.5 Corpus}
\label{chap:eras}
%==============================================================================

\emph{Commits: 2026-07-08 through 2026-07-13; the
\code{impl/era-pinned-helm-containers} branch, the \code{docker/era\_shim/}
package, and gotcha G13.}

\section{The problem: the measurement instrument itself has versions}

A large fraction of the audit corpus --- roughly \textbf{59\%}, the classic-track
\code{v0.2.4} and \code{v0.3.0} runs --- \emph{cannot} be replayed by the modern
runner, which pins \textsc{helm} 0.5.x and Python 3.11 and whose from-spec CLI
imports v0.5+ module paths that do not exist pre-v0.5. This is not a model problem;
it is a \emph{harness} problem. And it matters for reproducibility science because
the harness is the measurement instrument: the technical report had already shown
that the exact \code{pandas} version (2.0.x vs 2.2+) flips per-instance row
ordering on \code{entity\_matching}, i.e.\ the harness's own dependency versions
change the recorded instances.

The solution is to freeze the instrument at the era that produced each run: give
each era its own \emph{CPU-only} Docker image whose \textsc{helm} harness is
checked out at that era's release commit (v0.2.4 = \code{626d8609}, v0.3.0 =
\code{8ea285f7}), with era Python (3.10) and era dependency pins, while keeping
model inference \emph{out-of-process} on a modern vLLM lease. The container thus
holds only the scoring instrument fixed; the model server is shared with the
modern path.

\section{The design invariants}

\begin{itemize}[nosep]
  \item \textbf{One manifest = one era = one image = one measurement instrument.} A
        mixed-era source set is a hard error at make-manifest time
        (\code{eval\_audit/eras.py}). Era resolution keys on the path-derived
        \code{(public\_track, suite\_version)} signal.
  \item \textbf{Verbatim by-name replay.} A pre-v0.5 \code{adapter\_spec} has no
        \code{model\_deployment} field, so nothing is rewritten; routing to vLLM is
        purely by-name (the era shim registers a deployment under the exact official
        model name), and the materializer \emph{refuses} to insert a
        \code{model\_deployment} field into an era spec. Consequently
        \code{same\_deployment} resolves \code{unknown} for era pairs (both sides
        lack the field) --- the correct behavior, not a bug.
  \item \textbf{A standalone shim.} \code{docker/era\_shim/helm\_era\_shim/} never
        imports magnet or eval\_audit; it provides a flag-compatible from-spec CLI
        (\code{replay.py}) plus a backported OpenAI-compatible completions client, so
        the era image depends only on era \textsc{helm}.
  \item \textbf{An era$\leftrightarrow$image guard.} The image carries an
        \code{org.aiq.era} OCI label, checked against the manifest era at schedule
        time; a wrong pin fails loud.
\end{itemize}

\section{The v0.2.4-versus-v0.3.0 API skew}

A multi-angle review (verifying every era-API claim against the era \textsc{helm}
source via \code{git show <commit>:<path>}) found ten issues, the load-bearing ones
being cases where the shim had been written against v0.3.0 semantics and broke on
the older v0.2.4 API. Two are especially instructive:

\begin{itemize}[nosep]
  \item \textbf{Silent Together fallthrough (v0.2.4).} v0.2.4 lacks the
        register-deployments-from-path helper, so without an explicit
        \code{register\_model\_deployments\_from\_path} the deployments route to
        \code{TogetherClient} --- silently reintroducing the very hosted-API
        artifact the audit exists to eliminate.
  \item \textbf{pyhocon dot-splitting (both eras).} pyhocon path-splits credential
        keys on \code{.}, so a dotted model name like \code{pythia-6.9b} dies at the
        credential lookup. The fix writes a nested-key \code{credentials.conf} and
        puts \code{api\_key} in the client-spec args.
\end{itemize}

\section{A turnkey validation runbook and the 8~GB subject}

The era gates were rebuilt as \code{dev/era-tests/}, mirroring the \code{dev/e2e-
tests/} conventions, so one model runs end-to-end through both classic eras with
near-zero setup. The subject was swapped from \code{eleutherai/pythia-6.9b}
($\sim$14~GB, does not fit 8~GB) to \code{together/redpajama-incite-base-3b-v1},
the smallest corpus model with a \emph{full 74-run packet at both eras}
($\sim$2.8B params, $\sim$5.6~GB fp16). A ``validation ladder'' checks the
instrument in rungs: image sanity; instrument fidelity on a pandas-sensitive
\code{entity\_matching} run (byte-for-byte instance identity, no model); an
end-to-end single run; a full packet per era; and a per-scenario HF-fetch audit.
Several first-real-run bugs (a build-time import that named the wrong module for
the era; a \code{pandas}/\code{pyarrow} dependency-resolution skew; a fidelity rung
that diffed a \file{scenario\_state.json} the public corpus never ships --- switched
to the always-present \file{display\_requests.json}) reinforced the theme that
some bugs only surface at the first genuine build.

An important taxonomy point resurfaced here: rung 2 originally \emph{failed} on
scenarios whose data is no longer fetchable (e.g.\ MATH's script loader returns 401
on the 2026 Hub). That is an \emph{environment} filter, not an instrument failure,
so the rung was changed to classify data-fetch/auth crashes as SKIP --- the same
recipe-vs-reproducibility distinction, now applied to a historical-data gate.

\section{G13: a class path that resolves in no released \textsc{helm}}
\label{sec:g13}

The deepest archaeological finding of the internship concerns the original-paper
classic models (GPT-J 6B, GPT-NeoX 20B, OPT-66B) and is written up as gotcha G13
(\S\ref{app:gotchas}). Era-replaying these fails the shim's class preflight: their
\file{run\_spec.json} names \code{helm.benchmark.basic\_metrics.BasicMetric}, which
resolves in \emph{no} released \textsc{helm} version. An investigation over a
blob-less clone of the \textsc{helm} repository established that the stored path is
a once-migrated hybrid: a bulk \code{benchmark.}$\rightarrow$\code{helm.benchmark.}
prefix rewrite (at the \code{src/helm/} rename, 2022-11-16) applied to a run-spec
whose metric class was still \emph{flat} at production time. Reversing the naive
migration triangulates the producing code to a $\sim$4-week window
(2022-07-31 to 2022-08-26) of \emph{unreleased pre-v0.1.0 \textsc{helm}} --- after
scenarios were nested but before metrics were. The class-path \emph{shape} (flat
vs subpackage) is thus a sharper origin fingerprint than release dates. Because no
released version resolves these, and building the true origin instrument is
infeasible, the fix is a self-verifying declared class-path canonicalization in the
era shim: when a flat \code{helm.benchmark.*} class fails to resolve \emph{and} its
\code{helm.benchmark.metrics.*} relocation resolves to the same leaf, remap it on
the in-memory run-spec (recorded as a declared substitution). It fires only for the
known relocation --- a genuinely wrong era pin still fails loudly.

A related provenance fact was pinned along the way (memory
\emph{classic-officials-carried-forward}): the gptj/neox/opt/redpajama official
runs are \emph{byte-identical} across v0.2.4 and v0.3.0 --- carried forward, not
re-run --- so ``two eras agree'' is one public number measured by two instruments,
not two independent confirmations.

%==============================================================================
\chapter{New Model Families: GPT-OSS and Qwen}
\label{chap:newmodels}
%==============================================================================

\emph{Commits: 2026-07-11 through 2026-07-14; from-spec runbooks for
\code{openai/gpt-oss-20b} and a combined Qwen fan-out.}

With the from-spec, containerized, leased pipeline mature, two further model
families were brought in as reproduction subjects --- and as a test of whether the
machinery generalizes beyond OLMo.

\section{GPT-OSS-20B and the null-content crash}

The public GPT-OSS-20B rows were served via \code{together/gpt-oss-20b}, a chat
(harmony) client, so the faithful replay serves \emph{chat}, not the frozen older
preset's completions workaround. The design tension: GPT-OSS is a reasoning model,
and a chat turn can return \code{message.content=null} (reasoning-only,
\code{finish\_reason=length}), which un-patched \textsc{helm} crashes on. The root
cause was pinned precisely (memory \emph{gpt-oss-null-content-root-cause}): the
crash is a \emph{null-versus-empty-string serving difference} --- Together returns
\code{""}, vLLM returns \code{null} --- and normalizing
\code{null}$\rightarrow$\code{""} is faithful (\code{max\_tokens} and window were
ruled out as causes). The fix makes the null-safe chat clients normalize
\code{content:null}$\rightarrow$\code{""}, leaving \textsc{helm} untouched.
Discovery confirmed all four target run directories resolve 1:1 against the
84{,}966-run corpus with 0 ambiguities.

\finding{GPT-OSS-20B reproduces well.} The headline aggregate score-drift plot
(Figure~\ref{fig:gptoss}) shows local and public headline scores agreeing closely
across all four benchmarks --- \code{bbq} $0.967/0.963$, \code{gpqa} $0.594/0.610$,
\code{ifeval} $0.732/0.754$, \code{mmlu\_pro} $0.740/0.748$ (public/local) --- with
squared errors bounded by $\sim\!4\times10^{-4}$, roughly two orders of magnitude
tighter than the OLMo instruct drift on \code{ifeval}. This contrast is itself a
result: whether a family reproduces well is model-specific, and GPT-OSS-20B is a
positive case.

\begin{figure}[H]
\centering
\fbox{\parbox{0.92\textwidth}{\centering\small\itshape
[Figure: ``Headline Aggregate Score Drift'' heatmap for
\code{experiment\_name=gpt-oss-20b-from-spec}. One model row
(\code{openai/gpt-oss-20b}) $\times$ four benchmark columns (bbq, gpqa, ifeval,
mmlu\_pro), each cell annotated \code{P <public> / L <local>}, colored by
$(\text{local}-\text{public})^2$. Source: 07-14 status deck.]}}
\caption{GPT-OSS-20B reproduces the public headline metrics closely (squared error
$\le 4\times10^{-4}$). The rendered artifact is
\file{aggregate\_score\_diff\_headline.png} produced by
\code{build\_reports\_summary}.}
\label{fig:gptoss}
\end{figure}

\section{Qwen: generalizing the combined fan-out}

The eight public \textsc{helm} Qwen text models were reproduced as a single
\code{qwen-combined} multi-deployment from-spec fan-out --- the
\code{allenai-olmo-combined} analogue. The OLMo-specific ``combined preset''
helpers were generalized into \code{\_combined\_run\_entries} and
\code{\_build\_combined\_preset} so both families share one code path (safe to
generalize precisely because the OLMo combined tests pin the old output; regenerate
and diff the entry count).

Three decisions carried over as reusable rules:

\begin{itemize}[nosep]
  \item \textbf{Membership rule (the olmo-7b lesson).} Model size is \emph{never} a
        reason to exclude a member from the fan-out --- large models serialize under
        the lease while small ones co-host; the \emph{only} thing that forces a
        member split is discovery \emph{ambiguity} under the shared
        \code{precomputed\_root}. Propose all members and let the freeze preflight
        decide.
  \item \textbf{The whitelist is exactly reconstructible.} Filtering
        \code{official\_public\_index.csv} to classic-core + capabilities lands on
        the plan's per-model counts to the row ($85\times5 + 86 + 132\times2 = 775$),
        with two non-obvious exclusions (\code{banking77}, \code{bigcodebench}). All
        775 rows resolve with 0 ambiguities, so no member splits.
  \item \textbf{Protocol is answerable, not ``to confirm.''} \textsc{helm}'s own
        \code{model\_deployments.yaml} (read from the installed wheel) says base
        Qwen1.5 = \code{TogetherClient} (completions) and the chat/turbo models =
        \code{TogetherChatClient} (chat) --- so the serving protocol was resolved
        from the authoritative source, not guessed.
\end{itemize}

\section{The Qwen turbo window bug}
\label{sec:qwenwindow}

The last dated event of this chronology is a clean recurrence of the OLMo
chat-template-window interaction (Chapter~\ref{chap:olmo}), now on the Qwen turbo
endpoints. An \code{ifeval} smoke run on
\code{qwen/qwen2.5-72b-instruct-turbo} died with a
\code{ContextWindowExceededError}: the endpoint served \code{max\_model\_len:
4096}, but the two \code{*-instruct-turbo} endpoints are the only Qwen members that
replay the capabilities whitelist (\code{ifeval}, \code{mmlu\_pro}), whose official
\file{run\_spec.json} sets \code{max\_tokens=4096}. \textsc{helm}'s window service
truncates the \emph{prompt} but never clamps \code{max\_tokens}, so 4096 output
tokens fill the entire 4096 window and vLLM 400s with no room for a single prompt
token. The official Together serve ran a \code{max\_sequence\_length} of 128k, so
the overflow never surfaced upstream. The scoped fix raises those two endpoints to
\code{max\_model\_len: 8192} (4096 output + 4096 prompt headroom). This is a serving
configuration bug, not a reproducibility failure --- and it is a fitting last
example of the internship's recurring theme: the recipe under-specifies the
serving environment, and faithfully replaying the recipe surfaces exactly where
that under-specification bites.

%==============================================================================
\chapter{Reporting and Analysis Infrastructure}
\label{chap:reporting}
%==============================================================================

\emph{Commits throughout, concentrated 2026-07-08 through 2026-07-14. The reporting
layer is where reproducibility becomes legible, and much of the final fortnight was
spent making the aggregate story honest and comparable.}

\section{The pipeline the reports run on}

The analysis pipeline is read-only over audit results that already exist on disk
(``no model is run; no benchmark is downloaded''). Its stages are: (1) convert both
public and local runs to the \eeeshort artifact format
(\code{eval-audit-prepare-eee}); (2) compose a \emph{virtual experiment} from a
YAML manifest (\code{eval-audit-build-virtual-experiment}), applying a scope and
computing the coverage funnel; (3) per-packet core analysis
(\code{eval-audit-analyze-experiment}) through \code{NormalizedDiff}; (4) aggregate
publication (\code{build\_reports\_summary}). A \emph{virtual experiment} is the
organizing abstraction --- a declarative slice over existing audit data --- and the
\emph{packet} is the planner's canonical comparison unit (a bundle of normalized
components plus the comparisons to render between them).

\subsection{The three-level coverage funnel}
\label{sec:funnel}

The coverage funnel is the quantitative backbone of any ``fraction reproducible''
claim and reports three nested levels (Table~\ref{tab:funnel}). The reason all
three exist is gotcha G1: the raw \code{run\_spec\_hash} yields \emph{zero} matches
across \textsc{helm} releases even for semantically identical recipes, because
newer \textsc{helm} populates \code{adapter\_spec} fields older \textsc{helm} left
absent. The \emph{canonical} recipe hash collapses those known schema-evolution
fields before hashing; the gap between raw and canonical match is
\textsc{helm}-version churn, and the gap between canonical and logical match is
real recipe drift.

\begin{table}[H]
\centering
\small
\begin{tabularx}{\textwidth}{lY}
\toprule
\textbf{Level} & \textbf{Meaning} \\
\midrule
logical           & same scenario + model + augmentation (order-insensitive canonical key) \\
recipe-canonical  & + same scenario spec, prompt, decoding, \code{max\_train\_instances}, after schema-collapsing the run spec \\
recipe-identical  & byte-for-byte \code{run\_spec\_hash} match \\
\bottomrule
\end{tabularx}
\caption{The three-level coverage funnel. The report shows all three so that
\textsc{helm}-version churn (raw$\to$canonical) is separated from genuine recipe
drift (canonical$\to$logical).}
\label{tab:funnel}
\end{table}

\subsection{Sankey diagrams and the failure taxonomy}

Two Sankey diagrams render the funnel: Sankey~A (Universe$\to$Scope) narrows the
$13$k+ discovered runs to the in-scope set through ordered gates (structural,
metadata, open-weight, tag, deployment, size, manifest scope); Sankey~B
(Scope$\to$Reproduced$\to$Analyzed) carries in-scope rows through logical match,
recipe-canonical match, analyzed packet, and finally an agreement bucket. Crucially,
the failure-taxonomy colors and categories were made \emph{honest}: environment/
recipe exclusions and true reproducibility outcomes are distinct categories, so a
gated-model exclusion never reads as a reproducibility failure. This is the visual
realization of \S\ref{sec:distinction}.

\section{The aggregate score-drift heatmaps}

The final fortnight built the plots that summarize reproducibility at a glance ---
and, tellingly, most of the work was in making them \emph{correct and comparable}
rather than in drawing them.

\begin{itemize}
  \item \textbf{Per-metric aggregate-score-difference heatmap.} Instead of
        instance-level agreement, this shows, per core metric, the difference
        between the reproduced (local) and official (public) \emph{aggregate}
        scores; color is the signed difference on a symmetric diverging colormap
        (so $0$ is the white ``no drift'' midpoint), each cell annotated
        \code{P <public> / L <local>}. A deliberate data-source decision: read the
        signed means from the already-written run-level CSV
        (\code{core\_runlevel\_table.csv}) rather than re-serializing every packet's
        JSON --- it works on existing outputs with no re-render.
  \item \textbf{Headline holistic heatmap.} One plot showing every
        model$\times$benchmark pair, each benchmark using its \emph{headline}
        metric. \code{eval\_audit} had no primary-metric concept, but \textsc{helm}
        does (\code{run\_groups[].environment.main\_name} in the schema), so a
        curated \code{HEADLINE\_METRIC\_BY\_BENCHMARK} map (mmlu$\to$exact\_match,
        narrativeqa$\to$f1\_score, wmt\_14$\to$bleu\_4, \ldots) is used \emph{iff}
        that metric is actually present, falling back to a priority list otherwise;
        the plot names the actually-used metric on each axis tick, so the choice is
        transparent. The final version colors by squared error
        $(\text{local}-\text{public})^2$ (non-negative, emphasizes big drifts) and
        orients models-left / benchmarks-bottom to match the coverage matrix.
        Figures~\ref{fig:gptoss} and~\ref{fig:olmoheatmap} are instances.
\end{itemize}

\begin{figure}[H]
\centering
\fbox{\parbox{0.95\textwidth}{\centering\small\itshape
[Figure: ``Headline Aggregate Score Drift'' heatmap for
\code{experiment\_name=olmo-models-combined}. Six model rows (olmo-1.7-7b,
olmo-2-\{7b,13b,32b\}-instruct, olmo-7b, olmoe-1b-7b-instruct) $\times$ $\sim$12
benchmark columns. The base \code{olmo-7b} row reproduces the classic benchmarks
near-exactly (GSM $0.036/0.036$, MMLU $0.295/0.287$, narrative\_qa $0.597/0.595$,
wmt\_14 $0.097/0.098$); the instruct models show large drift on \code{ifeval}
(e.g.\ olmoe $0.628/0.731$; the darkest cells) and moderate drift on
\code{bbq}/\code{gpqa}/\code{mmlu\_pro}. The $\ddagger$ marks instance-level
fallback. Source: 07-14 status deck.]}}
\caption{The OLMo reproduction at a glance. Base-model classic benchmarks reproduce
near-exactly; instruct-model \code{ifeval} shows the largest drift --- consistent
with the chat-template-version (\code{add\_generation\_prompt}) and fp32-precision
findings of Chapter~\ref{chap:deploymatch} being most consequential on long-form,
instruct evaluations.}
\label{fig:olmoheatmap}
\end{figure}

\section{Correctness fixes that changed the numbers}

Several late fixes are worth recording because each one is a way the aggregate story
could have been \emph{silently wrong} --- and each is a lesson about honest
reporting:

\begin{itemize}[nosep]
  \item \textbf{Instance-level fallback for CoT benchmarks.} \code{ifeval},
        \code{gpqa}, and \code{mmlu\_pro} are chain-of-thought scenarios whose core
        metrics are \emph{instance-only} (\code{ifeval\_strict\_accuracy},
        \code{chain\_of\_thought\_correctness}); \textsc{helm} emits no run-level
        aggregate, so those benchmarks were silently \emph{absent} from the drift
        plot. The fix widens the data contract: fall back to the mean of the
        per-instance $0/1$ scores (which \emph{is} the accuracy), tagged with a
        $\ddagger$ so the provenance is visible. Widening the contract beats hacking
        the plot.
  \item \textbf{Metric-key granularity trap.} Run-level cells are keyed by full
        \textsc{helm} stat descriptions (``\code{f1\_score test on narrativeqa}'')
        while instance-fallback cells are keyed by bare ids
        (``\code{ifeval\_strict\_accuracy}''), so the curated headline map (bare
        names) silently failed to match run-level cells and every run-level
        benchmark fell to alphabetical-first. The fix normalizes to the metric
        \emph{family} at the matching boundary and selects the representative on the
        benchmark's \textsc{helm} \code{main\_split} (test vs valid differ, e.g.\
        narrative\_qa f1 $0.597$ test vs $0.661$ valid).
  \item \textbf{Stale-local dedupe.} An \code{olmo-7b}/mmlu cell showed a spurious
        $\sim\!0.15$ aggregate gap (public $0.295$ vs local $0.144$) despite
        near-exact instance agreement. The cause was not drift: every olmo-7b
        subject directory carried \emph{two} \code{official\_vs\_local} pairs --- a
        fresh run plus an un-demoted stale $0.0$ attempt --- so micro-averaging
        halved the local mean ($124$ rows instead of $62$). The fix restricts the
        accumulator to the canonical (first) comparison per report, with a warning
        per dropped attempt so a future genuine multi-local case surfaces instead of
        vanishing.
\end{itemize}

\section{The shared-gateway route registry}

A final infrastructure fix addressed a real operational failure: multiple runbooks
(olmo, qwen, gpt-oss, classic-together) share one standing \code{infer-stack} stack
but ship disjoint catalogs, and a converge under one catalog stripped another
runbook's still-live routes from the LiteLLM gateway --- presenting as a
\code{400 Invalid model name} and (because the minutes-long readiness wait silently
never passes once a route is stripped) as the operator's known ``failed to acquire
lease'' events. The fix persists an append-only \emph{route registry} in shared
state; every converge merges the invoking catalog with all live deployments and
renders the gateway route table from the \emph{whole} registry, so the rendered
config is a function of accumulated shared state rather than of which runbook
invoked converge --- byte-stable renders, so the gateway is never recreated. The
reusable insight: \emph{when a shared artifact is rebuilt by multiple writers,
render it from accumulated shared state, never from the invoking writer's
worldview.}

%==============================================================================
\chapter{Codebase Stewardship}
\label{chap:stewardship}
%==============================================================================

\emph{Commits: 2026-07-02, 2026-07-06, 2026-07-12. Three audit/simplification
passes interleaved with the research work.}

Reproducibility is a property of software, so the health of the codebase is part of
the scientific claim: a reviewer who cannot read the pipeline cannot trust its
verdicts. Three passes are worth recording as evidence of the ``code a researcher
can trust'' standard.

\begin{description}[style=nextline,leftmargin=1.25em]
  \item[Correctness audit (2026-07-02/03).] A full codebase audit produced a
        phased fix plan (P0/P1/P2 findings) implemented over $\sim$35 commits: e.g.\
        matching/joining on canonical keys rather than raw strings; correcting a
        perturbed/unperturbed agreement split; filtering non-finite scores at row
        construction; failure-taxonomy honesty (colors, categories, ordering);
        content-addressing generated \code{model\_deployments.yaml} filenames;
        tightening the HF-token permission window. Direct dependencies were declared
        and \code{transformers} pinned below 5 (a version whose \code{is\_offline\_
        mode} change breaks the \textsc{helm}$\to$\eeeshort rebuild).
  \item[Simplicity audit (2026-07-06).] An explicitly orthogonal ``reduce bloat''
        pass: $\sim$10{,}600 lines deleted (dead POC/oneoff trees, generated
        manifests, shims), \code{helm/diff.py} halved ($1{,}844\to844$),
        \code{PRESET\_CONFIGS} moved from a $1{,}154$-line dict to a YAML resource,
        the planning-doc set consolidated ($22\to11$ live docs), and the legacy
        agreement half of the comparison core finally retired (R-2), completing the
        migration onto \code{NormalizedDiff}. The methodological gem: pure-relocation
        refactors of plot-emitting code were proven correct by first
        \emph{characterizing the noise floor} (render \code{HEAD} twice, diff, and
        treat that file set as the only acceptable before/after diff --- necessary
        because plotly emits random div UUIDs).
  \item[Fourth simplification pass (2026-07-12).] A post-era ``finish what's in
        flight'' pass: Batches 1--3 landed in $22$ commits with the full
        \code{--run-slow} suite green ($638$ passed); \code{helm/} shrank
        $3{,}336\to\sim1{,}780$, \code{build\_reports\_summary.py} finished its split
        ($1{,}715\to\sim330$). A recurring review discipline shows through: audit-
        agent claims about \emph{why} code exists (``kept for tests'') were the ones
        to re-verify by grep --- liveness claims were reliable, purpose claims less
        so.
\end{description}

Two standing engineering disciplines recur across all three passes and are worth
lifting into the paper's reproducibility-methods note: (1) capture a
characterization baseline \emph{before} any behavior-changing refactor and gate on
\code{atol}$=10^{-9}$ equivalence; and (2) submodule hygiene --- land changes in the
submodule branch first, bump the gitlink as a deliberate separate commit, never
auto-commit a dirty gitlink.

%==============================================================================
\chapter{Thread B: Efficient Benchmarking}
\label{chap:effbench}
%==============================================================================

\emph{Repositories: \file{mrmr\_eval} and \file{dkps} (sibling repos, outside
\file{eval\_audit}); active $\sim$2026-06-01 through 2026-06-24; status decks 06-02,
06-09, 06-16, 06-23. Paused thereafter in favor of the HELM-reproduction thread.}

Running in parallel with the reproduction work for the first month was an
exploratory statistics thread on \emph{efficient benchmarking}: given that
evaluating a model on a full benchmark is expensive, can a small, well-chosen
\emph{coreset} of questions plus a regression predict the model's full-benchmark
score, and does a Data-Kernel Perspective-Space (\textsc{dkps}) representation help?
The experiment code lived in two sibling repositories forked from the relevant
prior work; the LoRA weight-space idea from the literature review
(Chapter~\ref{chap:litreview}) remained a proposal and produced no code.

\section{The two framings being compared}

\subsection{Feature selection + regression (Bowyer et al.)}

Following \emph{Efficient Benchmarking Is Just Feature Selection and Multiple
Regression}, the setup is a binary score matrix (models $\times$ questions). Each
question is a feature vector in $\mathbb{R}^M$ ($M$ source-model scores) and the
target is each model's overall accuracy. One \emph{selects} an informative coreset
of questions --- primarily with \textbf{mRMR} (Minimum-Redundancy-Maximum-
Relevance, in MIQ = relevance/redundancy and MID = relevance$-$redundancy variants)
--- and \emph{predicts} the full score with ridge or polynomial kernel-ridge
regression. The \code{benchpred} registry implements $\sim$100 method keys spanning
mRMR (ridge/kernel heads), IRT (\code{pirt}, \code{gpirt} = Gaussian-process IRT),
anchor points, Lasso, and random-sampling/search baselines, each optionally
``backfilled'' (append \code{+}) by re-fitting its coreset with ridge/kernel-ridge.

\subsection{DKPS: a model-centric, transductive predictor}

\textsc{dkps} (Data-Kernel Perspective-Space) is the model-centric alternative.
Given a pool of source models with known scores and a target model, all identified
only by their \emph{raw responses} to an aligned set of $N$ queries, the method:
(1) embeds every $(\text{model}, \text{query})$ response (via a text-embedding API,
or a one-hot encoding for multiple-choice tasks); (2) forms an inter-model
``data kernel'' --- a models$\times$models Euclidean distance matrix on the flattened
per-model response embeddings, normalized by $\sqrt{N}$; (3) computes the
\emph{perspective space} by classical MDS on that distance matrix (default 8
components); and (4) fits a linear regression from source coordinates to source
scalar scores and reads off the target's predicted score. It is \emph{transductive}
--- the target's responses participate in building the MDS space, so the space is
re-fit per prediction --- and it requires every model to answer the exact same
aligned queries under the same embedder.

\section{Experiment design}

The efficient-benchmarking sweeps concentrated on the HELM tasks that carry raw
responses (so \textsc{dkps} is applicable): \code{legalbench}, \code{math},
\code{med\_qa}, \code{wmt\_14}, plus \code{arc\_challenge}, \code{ifeval},
\code{commonsense}. The knobs swept were: number of source (reference) models
$\{20, 30, 40, 50\}$ (sometimes 60); coreset size $\{1,2,3,4,5,10,15\}\%$; three
seeds per cell. The reported metrics were MAE, RMSE, Kendall $\tau$, and Spearman
$\rho$ (Figure~\ref{fig:mrmr}).

\begin{figure}[H]
\centering
\fbox{\parbox{0.95\textwidth}{\centering\small\itshape
[Figure: \code{math} --- \code{mrmr\_eval} tutorial. A $2\times4$ grid of
RMSE(\%), MAE(\%), Kendall $\tau$, Spearman $\rho$; top row versus coreset size
(5/10/15\%), bottom row versus number of source models (20--50). Methods overlaid:
mRMR MIQ (k=5), gp-IRT, Search+, random\_sampling\_and\_learn, Random, AnchorPts,
Lasso. mRMR/kernel-anchor methods lead on RMSE/MAE; Lasso is the clear worst.
Source: 06-16 status deck.]}}
\caption{Representative efficient-benchmarking result on the \code{math} task.
Coreset-selection methods (mRMR, kernel anchor points) substantially outperform
Lasso and improve with coreset size; \textsc{dkps}-family predictors are competitive
only in specific regimes (many source models, tiny coresets).}
\label{fig:mrmr}
\end{figure}

\section{The DKPS+mRMR active-selection method}

The internship's own methodological contribution to this thread was an
\emph{adaptive} coreset selector, \code{dkps\_mrmr} (the realization of the
``kNN-mRMR online coreset selection'' proposal), building the coreset iteratively
and \emph{per target model}: seed with one global mRMR pick; then repeatedly
(a) \textsc{dkps}-embed source $\cup\ \{\text{target}\}$ on the current coreset,
(b) take the $k$ source models nearest the target in \textsc{dkps} space, and
(c) restrict the score matrix to those $k$ rows and take one greedy mRMR step to
append the next question; finally predict with a standard \textsc{dkps} transductive
regression. It subclasses both \code{MRMRPred} and \code{DKPSPred} to reuse the
mutual-information machinery and the embedding/regression, and required a new runner
capability so \code{predict} receives the full response matrix rather than a
pre-sliced coreset.

\section{Findings}

The concrete results (chiefly the \code{med\_qa} report) are candid, and their
honesty is itself worth preserving:

\begin{itemize}[nosep]
  \item \textbf{\textsc{dkps}-family methods are not general winners.} At $\sim$20
        source models and $\sim$50 queries, the feature-selection-paper methods
        (kernel anchor points, Search+, kernel-mRMR) outperform \textsc{dkps} and
        \code{dkps\_mrmr}; on the 40-source-model \code{med\_qa} row,
        \code{dkps\_mrmr} wins no coreset-size column.
  \item \textbf{\textsc{dkps} scales with source models where others do not.}
        \textsc{dkps} accuracy improves as the number of source models grows, which
        is not necessarily true of the other methods --- and this is where
        \code{dkps\_mrmr} shines: at a $1\%$ coreset its \code{med\_qa} MAE improves
        sharply from $9.06$ to $5.19$ as sources go $20\to50$, making it the
        \emph{single best method at 50 source models} in that corner.
  \item \textbf{\code{dkps\_mrmr} plateaus and is worst at large budgets.} It peaks
        early (MAE $\sim\!4.6\%$, $\rho\approx0.93$) and, unlike every other method,
        does not keep improving to $2.0$--$2.9\%$ MAE at $10$--$15\%$ coresets.
  \item \textbf{Cost and representation caveats.} Because the coreset is rebuilt per
        target, one trial costs $\sim$3~s (1\%) to $\sim$2.5~min (15\%), versus
        sub-second for plain \textsc{dkps}. mRMR is \emph{extremely slow} on
        continuous-score benchmarks (\code{wmt\_14}). And a one-hot response
        encoding (used when text embeddings are unavailable) exposes very limited
        information about a model compared to a text encoding --- for \code{med\_qa},
        with 5 choices but imperfect answer formatting, the one-hot dimension
        exceeds 5 and the representation is weak.
\end{itemize}

\section{What stayed a proposal}

Three lines from the ``directions'' slides were never implemented and are recorded
as open future work: \emph{variance-based early stopping} for adaptive
benchmarking; \emph{ensemble} predictors mixing feature-centric (AEA, Scales++) and
model-centric (\textsc{dkps}, Fluid) methods, e.g.\ re-weighting by \textsc{dkps}
distance from a reference model; and the entire \emph{LoRA weight-space} (W2T)
program (SVD canonicalization $\to$ transformer embedding $\to$ attribute
classification / performance prediction / adapter retrieval, and the exploratory
idea of learning a distribution over LoRA weights). This thread effectively
concluded at the end of June 2026 when attention returned fully to HELM
reproduction.

%==============================================================================
\chapter{Synthesis: What Is Reproducible, and Why}
\label{chap:synthesis}
%==============================================================================

\section{The headline scientific claims}

\emph{This section was revised on 2026-07-15 (see Appendix~\ref{app:revision}) to
scope its claims to the evidence actually in hand. The original draft asserted a
positive result --- ``a runnable subset of \textsc{helm} \emph{is} reproducible,
with residual drift that is small, concentrated, and attributable.'' That is a
plausible and motivating \emph{hypothesis}, but the experiments that would
establish it (a frozen population, freshly regenerated stores, end-to-end OLMo
confirmation on held-out instances, completed Qwen/classic runs, cross-machine
measurements, and a controlled ablation) are not all done. The claims below are the
defensible ones.}

\medskip
The more defensible thesis --- and the more scientifically interesting one --- is
that \emph{reproducing an \textsc{llm} benchmark is a layered system-identification
problem.} A public recipe is not a complete experimental record; the producing
measurement instrument spans the recipe, the model deployment, the software stack,
hardware-sensitive numerical behavior, and the artifact-processing history.
Controlled reconstruction can \emph{attribute} and sometimes \emph{close} apparent
reproducibility gaps; but when provenance has been lost or transformed, the original
experiment can be \emph{non-identifiable}. What the internship established, at the
evidence level actually reached:

\begin{enumerate}
  \item \textbf{The unrecorded execution substrate is real and consequential
        (established).} \textsc{helm} records the recipe and the model name but not
        the dtype, revision, tokenizer/template version, attention kernel, or
        software-stack versions (of 148 HF-client deployments, 19 pin dtype, 7 pin
        revision; Appendix~\ref{app:unrecorded}). Individual substrate parameters are
        \emph{identifiable to varying degrees}: the loader's dtype default and the
        deployment client class are readable from source/config; the tokenizer's
        special-token append is readable from the repo; the exact
        \code{transformers}/\code{torch}/CUDA versions and hardware are not.
  \item \textbf{Attribution pilot on OLMo (partly established, partly pending).} On a
        small ifeval sample, \code{float32} candidates best match the sampled public
        OLMo outputs, consistent with the deductive dtype-default argument
        (\S\ref{sec:fp32}); the OLMo-7B prompt-independent gibberish is attributable
        to a tokenizer EOS-append; OLMoE-instruct disagreement is attributable to
        \code{add\_generation\_prompt} chat-template drift. These are \emph{probe-} and
        \emph{mechanism-level} attributions on a pilot, \textbf{not} yet end-to-end
        \textsc{helm} reproductions confirmed on held-out instances (the ``confirm''
        step remains the top open experiment; \S\ref{chap:reporting} and the paper
        plan). ``Precision is \emph{the} dominant variable'' is a hypothesis
        supported within this pilot, not established across benchmarks or families.
  \item \textbf{A recoverable-vs-non-identifiable contrast (the strongest conceptual
        result).} OLMo is the \emph{recoverable} case: the missing provenance (dtype,
        tokenizer, template) can be reconstructed and its effect measured. The
        original-paper classic models (GPT-J, GPT-NeoX, OPT) are the
        \emph{non-identifiable} case: their run-specs name a class path that resolves
        in no released \textsc{helm} and triangulates to unreleased pre-v0.1.0 code
        (G13, \S\ref{sec:g13}), so the producing instrument cannot be uniquely
        identified or rebuilt. This contrast --- some lost provenance is
        experimentally recoverable, some is fundamentally underdetermined --- is a
        sharper contribution than a uniformly positive ``we closed the gaps'' story.
  \item \textbf{Non-runnable cases are a separate gate, not a reproducibility
        verdict.} Roughly half of the original broad grid could not be \emph{run at
        all} (gated datasets, revoked download access, missing credentials, packaging,
        resource contention). These are \emph{filtering/eligibility} reasons upstream
        of any disagreement measurement --- they answer ``can we run this?'', not
        ``did it reproduce?'' --- and must not be folded into a negative result.
\end{enumerate}

\section{The reusable methodology}

Independently of the specific numbers, the internship produced a reusable
methodology for reproducibility auditing that the paper can present as a
contribution:

\begin{itemize}[nosep]
  \item \textbf{Replay the resolved recipe, never reconstruct it}
        (Chapter~\ref{chap:fromspec}): from-spec replay of \file{run\_spec.json}
        eliminates a whole class of silent reconstruction drift.
  \item \textbf{Freeze the measurement instrument by era}
        (Chapter~\ref{chap:eras}): version-pinned harness containers make the
        scoring code a controlled variable, separable from the model.
  \item \textbf{Sweep, don't guess} (Chapter~\ref{chap:deploymatch}): for unrecorded
        execution parameters, an empirical sweep with feasibility-pruning (but not
        relevance-pruning) is the honest analogue of matching a recipe.
  \item \textbf{Classify every disagreement} (\S\ref{sec:distinction}): the
        recipe/environment-vs-reproducibility distinction, enforced through typed
        exclusion reasons, honest failure-taxonomy colors, and
        \code{comparability\_unknown:*} signals rather than silent fallbacks.
  \item \textbf{Report a distribution, not a point}: same-machine, cross-machine,
        and official-vs-local agreement, with per-metric tolerance sweeps, so that
        ``reproducible'' is quantified rather than asserted.
\end{itemize}

\section{The provenance recommendation for the field}

The clearest actionable output for the broader community is that \textsc{helm} (and
any benchmark harness) should record the \emph{execution substrate}, not just the
recipe and the model name. Three fields --- \code{model\_revision} (commit SHA),
\code{torch\_dtype}, and \code{transformers\_version} --- transitively close the
majority of the reproducibility gap, because the revision fixes the
tokenizer/template/generation-config state and the transformers version fixes the
precision default and the chat-template behavior. These belong in a machine-readable
provenance block (the existing \code{recipe\_facts} slot suffices with no schema
change), kept \emph{out} of the run-spec identity string so they do not break
official$\leftrightarrow$local pairing (Appendix~\ref{app:unrecorded}).

\section{Open threads at the close of this chronology}

Several items were scoped or built but not fully landed, and are the natural
next steps for the paper's experiments:

\begin{itemize}[nosep]
  \item The HuggingFace-in-process \emph{routing switch} --- the mechanism exists
        (reserve-only lease + fp32 \code{HuggingFaceClient}), but the default replay
        path does not yet auto-route a HuggingFaceClient official to it.
  \item GPU-host acceptance runs for several pipelines (era images, the qwen and
        gpt-oss fan-outs, the HF-in-process OLMoE exact-match test) were prepared but
        remained pending real GPU execution.
  \item Regenerating the stale on-disk OLMo report store so the corrected drift
        plots (instance-level fallback, metric-family resolver, stale-local dedupe)
        populate the published artifacts.
  \item Pinning model revisions through \code{eval\_audit}'s bundle boundary (a
        $\sim$3--4 file change) so generated bundles carry the SHA automatically.
\end{itemize}

%==============================================================================
\appendix
\chapter{The \textsc{helm} Gotchas Catalogue (G1--G13)}
\label{app:gotchas}
%==============================================================================

This appendix condenses the running ledger of undocumented or hard-to-see
\textsc{helm} behaviors accumulated while building the audit pipeline
(\file{docs/helm-gotchas.md}). Each is a concrete, reviewer-anticipating pitfall
that a reproducibility paper's appendix should address.

\begin{description}[style=nextline,leftmargin=1em]
  \item[\textbf{G1. \file{run\_spec.json} schema evolves silently between releases.}]
    Newer \textsc{helm} populates \code{adapter\_spec} fields older \textsc{helm}
    omitted (\code{chain\_of\_thought\_prefix}, \code{global\_suffix},
    \code{num\_trials}, \code{model\_deployment}, top-level \code{metric\_specs} /
    \code{groups} / \code{annotators}), so byte-for-byte hashing yields $0/N$
    matches across releases even for identical recipes. Workaround: a
    schema-collapsed \emph{canonical recipe hash} (\S\ref{sec:funnel}).
  \item[\textbf{G2. The local \code{suite} field is the experiment name, not a
    public-track version.}] A versioned join on \code{logical\_key + suite\_version}
    yields $0$ matches; the funnel detects a non-version-shaped local suite and
    reports the versioned join as not meaningful (\code{N/A}) rather than $0$.
  \item[\textbf{G3. The \code{huggingface/<model>} deployment alias does not always
    exist.}] Many open-weight runs were executed through the HF API, but
    \textsc{helm}'s registry ships deployment YAMLs for only a subset; a missing
    alias forces a different stack (vLLM/etc.), introducing deployment-substitution
    drift. Workaround: register a custom local deployment.
  \item[\textbf{G4. \code{model\_deployment} semantics differ across versions.}]
    Older \textsc{helm} did not record the field (defaulting to ``the HF API'');
    newer does. A local \code{huggingface/<model>} and an official
    \code{<MISSING>} are semantically identical but hash differently ---
    schema-evolution, dropped from the canonical hash.
  \item[\textbf{G5. \code{adapter\_spec.instructions} differs between releases on
    identical scenarios.}] MMLU official prepends ``Answer with only a single
    letter.''; several LegalBench scenarios prepend a list-the-allowed-labels
    preamble. This is \emph{genuine} recipe drift (a different prompt), surfaced via
    the \code{same\_instructions} fact --- not to be silently normalized. (The BBQ
    \code{output\_format\_instructions=mcqa} drift of \S\ref{sec:canonkey} is the
    same species.)
  \item[\textbf{G6. The \code{classic} corpus moved bucket prefixes.}] Classic runs
    migrated from the bucket root to \code{classic/benchmark\_output/runs/<ver>/},
    mirroring every other suite; a special-case in the downloader was removed.
  \item[\textbf{G7. Same weights $\ne$ same output across serving stacks.}] Even
    greedy, HF \code{generate()} and vLLM differ in kernels, tokenizer batching,
    stop-sequence handling, unspecified sampling defaults, and EOS/pad handling ---
    a $3$--$6\%$ per-instance flip rate on multiple-choice, $\sim$30--40\% on
    long-form. Mitigations: match the original client, pin decoding, or document as
    \code{deployment\_drift}. (Chapter~\ref{chap:deploymatch} is the deep dive.)
  \item[\textbf{G8. \code{metric\_specs} schema evolution.}] Top-level
    \code{metric\_specs} differ on the majority of packets from restructuring, not a
    recipe change; excluded from the canonical hash --- compare the metric
    \emph{output}, not the declaration.
  \item[\textbf{G9. Local index \code{model\_deployment} is audit-specific.}] Custom
    local deployment names (e.g.\ \code{kubeai/vicuna-...-no-chat-template-local})
    do not appear in the public run spec; the logical-key join collapses across
    deployments and \code{same\_deployment} correctly reports \code{no}.
  \item[\textbf{G10. Public-track \code{suite\_version} is not the \textsc{helm}
    release version.}] The same \file{run\_spec.json} can appear identically in
    \code{v0.2.4} and \code{v0.3.0}; use \code{run\_spec\_hash} to detect
    recipe-identical duplicates across suites. (The era-replay containers key the
    \emph{instrument} on \code{(public\_track, suite\_version)}, which is the right
    granularity for the classic track --- see Chapter~\ref{chap:eras}.)
  \item[\textbf{G11. \file{per\_instance\_stats.json} corruption on giant runs.}]
    Certain msmarco runs were truncated mid-write during upload (consistent
    byte-offset \code{JSONDecodeError}); recently re-uploaded versions parse cleanly.
    Redownload on a size mismatch.
  \item[\textbf{G12. \code{run\_spec\_hash} canonicalization differs by release.}]
    \textsc{helm}'s own canonicalization is version-dependent, so semantically
    identical specs hash differently across releases; the coverage-side canonical
    hash applies a stable normalization on top.
  \item[\textbf{G13. Classic-corpus \code{class\_name} paths resolve in \emph{no}
    \textsc{helm} version.}] The stored \code{helm.benchmark.basic\_metrics.BasicMetric}
    (helm-prefix + flat) is a once-migrated hybrid --- a naive
    \code{benchmark.}$\to$\code{helm.benchmark.} rewrite of a still-flat production
    path --- pinning the producing code to unreleased pre-v0.1.0 \textsc{helm}
    (2022-07-31 to 08-26). Class-path \emph{shape} is the origin fingerprint; fixed
    by a self-verifying declared class-path canonicalization in the era shim
    (\S\ref{sec:g13}).
\end{description}

%==============================================================================
\chapter{The Unrecorded-Parameter Taxonomy}
\label{app:unrecorded}
%==============================================================================

Condensed from \file{docs/helm-unrecorded-deployment-params.md}. \textsc{helm}
serializes the recipe and the model's name but almost nothing about \emph{how} the
numbers were computed. Empirically, of 148 HuggingFaceClient deployments in
\textsc{helm}'s \code{model\_deployments.yaml}, only \textbf{19 pin
\code{torch\_dtype}} (all \code{bfloat16}) and only \textbf{7 pin a model
\code{revision}}. The parameters, ranked by effect on a greedy (temperature-0)
output:

\begin{longtable}{@{}p{0.02\textwidth}p{0.30\textwidth}p{0.60\textwidth}@{}}
\toprule
\textbf{Tier} & \textbf{Parameter} & \textbf{Why it matters} \\
\midrule
\endhead
1 & Model weight revision (SHA) & Same name $\ne$ same weights; authors re-upload
    checkpoints/configs/tokenizers. The single highest-leverage field --- fixes the
    tokenizer, chat template, and generation config as-of-then. \\
1 & Load precision / dtype & Unpinned $\Rightarrow$ transformers-version-dependent
    default (pre-v5 = float32). \emph{Confirmed decisive}: HF fp32 reproduces OLMoE
    exactly (\S\ref{sec:fp32}). \\
1 & Quantization & GPTQ/AWQ/fp8/bnb change the weights outright. \\
1 & Software-stack versions & The meta-parameter: transformers sets the fp32-vs-bf16
    default \emph{and} whether the chat template honors \code{add\_generation\_prompt};
    engine/flash-attn versions change kernels. Nowhere in the run dir. \\
\midrule
2 & Attention implementation & eager/sdpa/flash/PagedAttention change reduction order
    $\Rightarrow$ different logits at temp 0. \\
2 & Matmul / TF32 precision & For fp32 especially, true-fp32 vs TF32 matmuls differ. \\
2 & Hardware / GPU model & Determines kernels, TF32 behavior, whether fp32 MoE fits. \\
2 & Device topology / TP & Multi-GPU all-reduce order differs from single-GPU. \\
2 & Batch composition & Kernels are not batch-invariant; same prompt, different
    batch neighbors, different logits. \\
\midrule
3 & Tokenizer + chat-template version & The \code{add\_generation\_prompt} drift
    (\S\ref{sec:addgenprompt}); the templated prompt is never stored. \\
3 & \code{add\_special\_tokens} / BOS & The OLMo-7B EOS-append case
    (\S\ref{sec:eosappend}). \\
3 & \code{generation\_config.json} defaults & Injects unspecified fields
    (repetition penalty, eos/pad ids). \\
3 & Stop-sequence / detokenization & HF stops on token ids, vLLM on string match
    $\Rightarrow$ different truncation and scored text. \\
\midrule
4 & RNG seed; KV-cache dtype & Mostly for sampling (temperature $>0$) and
    serving-engine reproductions. \\
\bottomrule
\end{longtable}

\noindent\textbf{The one non-obvious lever.} Several Tier-3 gaps collapse into
Tier-1: the tokenizer, chat template, and generation config all live \emph{inside}
the model repo, so pinning the model revision recovers all of them as-of-then.
\code{model\_revision} + \code{torch\_dtype} + \code{transformers\_version} close
the majority of the gap. Both weight-loading engines already accept \code{revision}
as a data-only edit (infer-stack \code{catalog.yaml}; \textsc{helm}'s
\code{client\_spec.args}); the only code gap is that \code{eval\_audit}'s bundle
boundary drops it (a $\sim$3--4 file fix). These fields must live in the
\code{recipe\_facts}/provenance block, \emph{not} the run-spec identity string, or a
SHA in the pairing key would break official$\leftrightarrow$local pairing.

%==============================================================================
\chapter{Deliverables and Timeline}
\label{app:deliverables}
%==============================================================================

\section{Internship timeline at a glance}

\begin{longtable}{@{}p{0.16\textwidth}p{0.80\textwidth}@{}}
\toprule
\textbf{Dates (2026)} & \textbf{Work} \\
\midrule
\endhead
May 15 -- Jun 2  & Onboarding; literature review (LoRA/W2T, efficient benchmarking).
    Thread B begins. \\
Jun 3 -- Jun 4   & First commits: e2e adapter config and tests (Chapter~\ref{chap:e2e}). \\
Jun 9            & Corpus cartography / filter statistics (Chapter~\ref{chap:cartography}). \\
Jun 11 -- Jun 12 & Repo refactor (Phases 0--2) and the unified comparison core
    (Phase 3 / \code{NormalizedDiff}, judge registry) (Chapter~\ref{chap:core}). \\
Jun 12 -- Jun 18 & OLMo campaign begins: tokenizers, chat-template window budget,
    the 114-packet canonical-key fix (Chapter~\ref{chap:olmo});
    Thread~B \code{dkps\_mrmr} lands (Jun 24). \\
Jun 16 -- Jun 24 & Containerized HELM execution and the infer-stack leasing pipeline
    (Chapter~\ref{chap:containers}). \\
Jun 25 -- Jun 30 & From-spec replay: run-spec verbatim replay, deployment rewrite,
    rel-path addressing, content-addressed caching (Chapter~\ref{chap:fromspec}). \\
Jun 30 -- Jul 10 & Deployment matching: OLMo-7B EOS-append, the deployment-match
    tool, \code{add\_generation\_prompt} drift, the float32 discovery, the
    unrecorded-parameter taxonomy (Chapter~\ref{chap:deploymatch}). \\
Jul 2 / 6 / 12   & Three codebase audit / simplification passes
    (Chapter~\ref{chap:stewardship}). \\
Jul 8 -- Jul 13  & Era-pinned containers, the shim, redpajama-3b runbook, classic
    Together models, the G13 archaeology (Chapter~\ref{chap:eras}). \\
Jul 11 -- Jul 14 & GPT-OSS-20B and Qwen fan-outs (Chapter~\ref{chap:newmodels});
    aggregate score-drift reporting fixes, route registry
    (Chapter~\ref{chap:reporting}). \\
\bottomrule
\end{longtable}

By commit count: 452 authored commits, distributed as 181 in June and 271 in July;
the busiest days were 2026-07-06 (50 commits, the simplicity audit) and 2026-07-12
(39, the fourth simplification pass).

\section{Key artifacts produced}

\paragraph{Reproduction runbooks} (\file{reproduce/}, \file{dev/}):
\code{olmo\_models\_combined}, \code{gpt\_oss\_20b\_from\_spec},
\code{qwen\_models\_combined}, \code{classic\_together\_combined},
\code{dev/e2e-tests} (phi-2), \code{dev/era-tests} (redpajama-3b).

\paragraph{Tools and libraries}: \code{dev/tools/deployment\_match/} (the
serve-and-probe deployment matcher with \code{hf-probe}, \code{hf-match},
fp32-TP, and \code{compare\_prompt.py}); \code{docker/era\_shim/} (the standalone
era from-spec replay CLI + OpenAI-compatible client); \code{eval\_audit/eras.py}
and \code{docker/eras.yaml} (era registry); \code{eval\_audit/hf\_inprocess.py}
(reserve-GPU in-process reproduction); \code{eval\_audit/normalized/diff.py}
(\code{NormalizedDiff}, the unified comparison core);
\code{eval\_audit/judge\_registry.py} and \code{eval\_audit/metrics\_taxonomy.py}.

\paragraph{Documentation} (the intended paper appendix material):
\file{helm-gotchas.md} (G1--G13), \file{vllm-vs-huggingface-deployment-match.md},
\file{helm-unrecorded-deployment-params.md}, \file{eee-vs-helm-metadata.md},
\file{pipeline.md}, \file{architecture.md}, \file{container-execution.md}, and the
19 planning documents under \file{docs/planning/}.

\paragraph{Thread~B} (sibling repos): \file{mrmr\_eval} (the \code{benchpred}
efficient-benchmarking harness with the \code{dkps}/\code{dkps\_mrmr} methods and the
\code{med\_qa} report) and \file{dkps} (the vendored DKPS library and HELM example
runners).

\section{A note on method and provenance of this document}

This chronology was reconstructed after the fact from the primary sources listed in
Chapter~\ref{chap:framing}. Where numbers are quoted (agreement ratios, token
counts, corpus statistics, the fp32 sweep table, the efficient-benchmarking results)
they are taken directly from the status decks, the developer journals, or the
committed reports; where a claim is a design decision or a lesson, it is drawn from
the corresponding journal entry or planning document. The development was carried
out with heavy use of an AI pair-programming harness (recorded in the first-person
developer journals); the research direction, design decisions, GPU-host execution,
and interpretation are the author's. Readers preparing the academic manuscript
should treat the \S-cross-referenced figures and tables here as pointers into those
primary artifacts, which remain the ground truth.

%==============================================================================
\chapter{Interpretive Revision (2026-07-15): Epistemic Status of the Central Claims}
\label{app:revision}
%==============================================================================

The body of this chronology reconstructs, in a single after-the-fact narrative
voice, what was believed and reported across the internship. Some of those claims
--- especially around the \code{float32} result --- were stated more strongly than
the evidence currently supports. Rather than rewrite the dated narrative (which
would erase what was believed at each point), this appendix records the corrected
epistemic status in one place. It was prompted by an external review of the
chronology and the paper plan, and its conclusions are folded into
[`docs/planning/tmlr-paper-plan.md`](../planning/tmlr-paper-plan.md).

\section{A workflow-status vocabulary}

Claims in the body should be read against \emph{where in the pipeline} their
evidence sits. These are not interchangeable, and the body sometimes chose the
strongest available interpretation:

\begin{enumerate}[nosep]
  \item \textbf{Infrastructure implemented} --- code exists, unit-tested on a CPU host.
  \item \textbf{Smoke-tested execution} --- ran end-to-end on a tiny input, once.
  \item \textbf{Deployment-match probe} --- an oracle-scored config sweep over a small
        sample, possibly using non-default request knobs. \emph{Discovery, not
        confirmation.}
  \item \textbf{Exact-spec run} --- the official \file{run\_spec.json} executed
        verbatim on the real path.
  \item \textbf{Complete \textsc{helm} artifact generation} --- normal
        \file{scenario\_state.json}/\file{stats.json} produced.
  \item \textbf{Final comparison report} --- run through the standard pair-report
        pipeline.
  \item \textbf{Fresh vs.\ stale store} --- regenerated with current code vs.\ a store
        predating a relevant fix.
\end{enumerate}

The single most important correction: the \code{float32} result (\S\ref{sec:fp32})
sits at level~3 (probe / discovery), not levels~4--6. The ``confirm'' step that
would move it to a reproduction result --- land the winning config on the ordinary
\textsc{helm} path, run the exact public \file{run\_spec.json}, generate normal
artifacts, compare via pair-report, and evaluate on \emph{held-out} instances --- has
not been run for any model. Configuration discovery and held-out confirmation must be
separated, or a config fit to a small sample is reported as if it were an identified
fact.

\section{A finer disagreement taxonomy}

\S\ref{sec:distinction}'s binary --- ``recipe/environment failure'' vs.\ ``true
reproducibility failure'' --- is too coarse, and the phrase ``true reproducibility
failure'' should be \emph{avoided} whenever execution provenance is incomplete. Keep
the upstream run-vs-not-run gate (non-runnable = a filtering/eligibility reason, not
a disagreement), and for observed disagreements use:

\begin{enumerate}[nosep]
  \item \textbf{Recipe mismatch} --- the prompts/decoding/scenario differ (e.g.\ the
        BBQ \code{output\_format\_instructions} drift, G5).
  \item \textbf{Deployment mismatch} --- serving engine / client class differs (vLLM
        vs.\ hosted API; \code{HuggingFaceClient} vs.\ vLLM).
  \item \textbf{Execution-instrument mismatch} --- precision, tokenizer/template
        version, attention kernel, device placement, software-stack versions.
  \item \textbf{Artifact-interpretation or migration mismatch} --- the pandas
        row-order / instance-id case (EEE Case~A); the G13 class-path migration.
  \item \textbf{Residual disagreement under controlled equivalence} --- the only
        category that resembles conventional repeatability failure: everything
        controllable is matched and a gap remains.
  \item \textbf{Historically non-identifiable reproduction} --- the surviving
        artifacts and releases do not uniquely specify what faithful reproduction
        would even mean (GPT-J / GPT-NeoX / OPT via G13). Qualitatively different from
        (5): not ``it failed to repeat'' but ``the target is underdetermined.''
\end{enumerate}

Most of what the body called ``deployment-boundary artifacts'' are categories
(2)--(4) --- recoverable by controlling the substrate. Category (6) is the negative
counterpart and is a first-class result, not a gap to be closed.

\section{Corrected-claims table}

\begin{longtable}{@{}p{0.42\textwidth}p{0.54\textwidth}@{}}
\toprule
\textbf{As stated in the body} & \textbf{Corrected status} \\
\midrule
\endhead
``The official OLMoE run --- and every OLMo HF-client run --- executed in
\code{float32}.'' & The \emph{loader code path} defaults to fp32 under every
\code{transformers} version consistent with the run date (a deductive claim, the
stronger leg), \emph{conditional} on the unrecorded \code{transformers}$<$5. Likely,
not proven. \\
``fp32 reproduces the official completions \emph{exactly}.'' & Among tested configs,
fp32 candidates \emph{best matched the sampled} public outputs ($n\approx12$ ifeval),
using a probe-only \code{add\_special\_tokens=False} knob. Not an end-to-end HELM
reproduction; not held-out. \\
``all four officials reproduce at fp32.'' & Probe-level, ifeval-only, instruct-only.
OLMoE via HF fp32; dense OLMo-2 via vLLM fp32; the HF probe does \emph{not} reproduce
OLMo-2 (an open execution-sensitivity finding, below). \\
``Precision is the dominant hidden variable.'' & Supported \emph{within} the OLMo
ifeval pilot; not established across benchmarks or model families. \\
``the OLMo-2 HF divergence is the probe mis-generating, not science.'' & Reframed: a
\emph{current} HF stack failing to reproduce an \emph{apparently} HF-produced
historical result under byte-identical prompt tokens is itself an unresolved
systematic divergence under prompt-token equivalence --- evidence for the thesis. \\
``residual drift is small, concentrated, and almost always attributable.'' & A
motivating \emph{hypothesis}, pending the confirm step, fresh stores, cross-machine
data, and a controlled ablation. Not yet established. \\
OLMo/GPT-OSS heatmaps ``exist''. & Report stores exist (Jul 14) but carry
stale-local numbers on disk; must be regenerated with current code before citing. \\
\bottomrule
\end{longtable}

None of these corrections weaken the \emph{methodological} contributions (from-spec
replay, era-pinned instruments, sweep-don't-guess attribution, the provenance
taxonomy) or the \emph{conceptual} contribution (recoverable vs.\ non-identifiable
provenance). They scope the \emph{empirical} claims to what has actually been run.

\end{document}
